\begingroup\tcbset{breakable/.default=false}
\renewcommand{\Needspace}[1]{\par\begingroup
\dimen0=#1\relax\vskip0pt plus\dimen0\penalty-100
\vskip0pt plus-\dimen0\vskip\dimen0\penalty9999\vskip-\dimen0
\vskip0pt\endgroup}
\chapter{成像与光学仪器}


\graphicspath{{/Users/dongxingxi/Documents/物理竞赛一轮资料/sample/xes-original/forty-third-section/}{/Users/dongxingxi/Documents/物理竞赛一轮资料/sample/xes-original/forty-third-section/assets/}}


\section{成像概念、光学器件与近轴条件}


\newcommand{\SphereRefractionPanel}{%
\begin{center}\begin{minipage}{\linewidth}\centering
\begin{tikzpicture}[>=Latex,font=\small,every node/.style={fill=none}]
\draw[thick] plot coordinates {(0.3616959,1.1471529) (0.3478034,1.1270521) (0.3341575,1.1067831) (0.3207603,1.0863489) (0.3076137,1.0657526) (0.2947197,1.0449971) (0.2820802,1.0240857) (0.2696972,1.0030215) (0.2575724,0.9818075) (0.2457077,0.9604470) (0.2341048,0.9389431) (0.2227655,0.9172991) (0.2116915,0.8955182) (0.2008844,0.8736036) (0.1903459,0.8515586) (0.1800775,0.8293865) (0.1700807,0.8070906) (0.1603570,0.7846742) (0.1509079,0.7621408) (0.1417349,0.7394935) (0.1328391,0.7167359) (0.1242221,0.6938713) (0.1158851,0.6709031) (0.1078293,0.6478348) (0.1000559,0.6246698) (0.0925661,0.6014116) (0.0853610,0.5780636) (0.0784417,0.5546293) (0.0718092,0.5311122) (0.0654645,0.5075159) (0.0594085,0.4838438) (0.0536423,0.4600995) (0.0481665,0.4362865) (0.0429820,0.4124084) (0.0380897,0.3884687) (0.0334902,0.3644711) (0.0291842,0.3404190) (0.0251724,0.3163161) (0.0214553,0.2921661) (0.0180336,0.2679724) (0.0149077,0.2437387) (0.0120781,0.2194686) (0.0095452,0.1951658) (0.0073094,0.1708338) (0.0053710,0.1464764) (0.0037304,0.1220971) (0.0023877,0.0976995) (0.0013432,0.0732874) (0.0005970,0.0488644) (0.0001493,0.0244340) (0.0000000,0.0000000) (0.0001493,-0.0244340) (0.0005970,-0.0488644) (0.0013432,-0.0732874) (0.0023877,-0.0976995) (0.0037304,-0.1220971) (0.0053710,-0.1464764) (0.0073094,-0.1708338) (0.0095452,-0.1951658) (0.0120781,-0.2194686) (0.0149077,-0.2437387) (0.0180336,-0.2679724) (0.0214553,-0.2921661) (0.0251724,-0.3163161) (0.0291842,-0.3404190) (0.0334902,-0.3644711) (0.0380897,-0.3884687) (0.0429820,-0.4124084) (0.0481665,-0.4362865) (0.0536423,-0.4600995) (0.0594085,-0.4838438) (0.0654645,-0.5075159) (0.0718092,-0.5311122) (0.0784417,-0.5546293) (0.0853610,-0.5780636) (0.0925661,-0.6014116) (0.1000559,-0.6246698) (0.1078293,-0.6478348) (0.1158851,-0.6709031) (0.1242221,-0.6938713) (0.1328391,-0.7167359) (0.1417349,-0.7394935) (0.1509079,-0.7621408) (0.1603570,-0.7846742) (0.1700807,-0.8070906) (0.1800775,-0.8293865) (0.1903459,-0.8515586) (0.2008844,-0.8736036) (0.2116915,-0.8955182) (0.2227655,-0.9172991) (0.2341048,-0.9389431) (0.2457077,-0.9604470) (0.2575724,-0.9818075) (0.2696972,-1.0030215) (0.2820802,-1.0240857) (0.2947197,-1.0449971) (0.3076137,-1.0657526) (0.3207603,-1.0863489) (0.3341575,-1.1067831) (0.3478034,-1.1270521) (0.3616959,-1.1471529)};
\draw[gray,thin,->] (-4.6000000,0.0000000)--(7.8995696,0.0000000);
\fill (-4.0000000,0.0000000) circle (1.2pt);
\node[below,fill=none,inner sep=2pt] at (-4.0000000,0.0000000) {$S$};
\fill (0.0000000,0.0000000) circle (1.2pt);
\node[below left,fill=none,inner sep=2pt] at (0.0000000,0.0000000) {$V$};
\fill (2.0000000,0.0000000) circle (1.2pt);
\node[below,fill=none,inner sep=2pt] at (2.0000000,0.0000000) {$C$};
\fill (7.4995696,0.0000000) circle (1.2pt);
\node[below,fill=none,inner sep=2pt] at (7.4995696,0.0000000) {$P$};
\fill (0.0482051,0.4364592) circle (1.2pt);
\node[above left,fill=none,inner sep=2pt] at (0.0482051,0.4364592) {$A$};
\draw[gray,dashed,thin] (2.0000000,0.0000000)--(0.0482051,0.4364592);
\draw[DeepBlue,thick,->] (-4.0000000,0.0000000)--(0.0482051,0.4364592);
\draw[DeepBlue,thick,->] (0.0482051,0.4364592)--(7.4995696,0.0000000);
\node[above,fill=none,inner sep=2pt] at (-2.0000000,1.0000000) {$n_1=1$};
\node[above,fill=none,inner sep=2pt] at (3.0000000,1.0000000) {$n_2=2$};
\end{tikzpicture}
\figcap{球面顶点 $V$、球心 $C$ 与球面点 $A$；图中 $P$ 为所示有限孔径光线的轴交点。}\label{ch28:fig:s43-20-1}
\end{minipage}\end{center}
}

\newcommand{\VirtualSignPanel}{%
\begin{center}\begin{minipage}{\linewidth}\centering
\begin{tikzpicture}[>=Latex,font=\small,every node/.style={fill=none}]
\begin{scope}[shift={(-3.2000000,0.0000000)}]
\draw[thick] (0.0000000,-1.4000000)--(0.0000000,1.4000000);
\draw[gray,thin,->] (-2.7000000,0.0000000)--(2.7000000,0.0000000);
\draw[DeepBlue,thick,->] (-2.0000000,0.0000000)--(0.0000000,-0.5500000);
\draw[DeepBlue,thick,->] (0.0000000,-0.5500000)--(2.4606300,-0.9919275);
\draw[gray,dashed,thin] (-3.0623724,0.0000000)--(0.0000000,-0.5500000);
\draw[DeepBlue,thick,->] (-2.0000000,0.0000000)--(0.0000000,0.5500000);
\draw[DeepBlue,thick,->] (0.0000000,0.5500000)--(2.4606300,0.9919275);
\draw[gray,dashed,thin] (-3.0623724,0.0000000)--(0.0000000,0.5500000);
\fill (-3.0623724,0.0000000) circle (1.2pt);
\node[below,fill=none,inner sep=2pt] at (-3.0623724,0.0000000) {};
\draw[gray,dashed,thin] (-3.0623724,0.0000000)--(-3.0623724,-0.6000000);
\node[below,fill=none,inner sep=2pt] at (-3.0623724,-0.6500000) {虚像};
\fill (-2.0000000,0.0000000) circle (1.2pt);
\node[below,fill=none,inner sep=2pt] at (-2.0000000,0.0000000) {};
\draw[gray,dashed,thin] (-2.0000000,0.0000000)--(-2.0000000,0.8500000);
\node[above,fill=none,inner sep=2pt] at (-2.0000000,0.9000000) {实物};
\node[below,fill=none,inner sep=2pt] at (0.0000000,-1.6000000) {空气到玻璃};
\end{scope}
\begin{scope}[shift={(3.2000000,0.0000000)}]
\draw[thick] (0.0000000,-1.4000000)--(0.0000000,1.4000000);
\draw[gray,thin,->] (-2.7000000,0.0000000)--(2.7000000,0.0000000);
\draw[DeepBlue,thick,->] (-2.4105135,-1.2128912)--(0.0000000,-0.5500000);
\draw[gray,dashed,thin] (0.0000000,-0.5500000)--(2.0000000,0.0000000);
\draw[DeepBlue,thick,->] (0.0000000,-0.5500000)--(1.2687483,0.0000000);
\draw[DeepBlue,thick,->] (-2.4105135,1.2128912)--(0.0000000,0.5500000);
\draw[gray,dashed,thin] (0.0000000,0.5500000)--(2.0000000,0.0000000);
\draw[DeepBlue,thick,->] (0.0000000,0.5500000)--(1.2687483,0.0000000);
\fill (1.2687483,0.0000000) circle (1.2pt);
\node[below,fill=none,inner sep=2pt] at (1.2687483,0.0000000) {};
\draw[gray,dashed,thin] (1.2687483,0.0000000)--(1.2687483,-0.6000000);
\node[below,fill=none,inner sep=2pt] at (1.2687483,-0.6500000) {实像};
\fill (2.0000000,0.0000000) circle (1.2pt);
\node[below,fill=none,inner sep=2pt] at (2.0000000,0.0000000) {};
\draw[gray,dashed,thin] (2.0000000,0.0000000)--(2.0000000,0.8500000);
\node[above,fill=none,inner sep=2pt] at (2.0000000,0.9000000) {虚物};
\node[below,fill=none,inner sep=2pt] at (0.0000000,-1.6000000) {玻璃到空气};
\end{scope}
\end{tikzpicture}
\figcap{左图出射光的反向延长线交于虚像；右图虚线为未遇界面时的会聚路径。有限孔径轴交点随角度略变。}\label{ch28:fig:s43-20-2}
\end{minipage}\end{center}
}

\newcommand{\EyePanel}{%
\begin{center}\begin{minipage}{\linewidth}\centering
\begin{tikzpicture}[>=Latex,font=\small,every node/.style={fill=none}]
\draw[thick] (0,0) ellipse (2.4 and 1.3);
\draw[thick] (-1.3,-.7) .. controls (-1.65,-.2) and (-1.65,.2) .. (-1.3,.7) .. controls (-.95,.2) and (-.95,-.2) .. cycle;
\draw[DeepBlue,thick,->] (-3.6000000,-0.5500000)--(-1.3000000,-0.5500000);
\draw[DeepBlue,thick,->] (-1.3000000,-0.5500000)--(2.4000000,0.0000000);
\draw[DeepBlue,thick,->] (-3.6000000,0.5500000)--(-1.3000000,0.5500000);
\draw[DeepBlue,thick,->] (-1.3000000,0.5500000)--(2.4000000,0.0000000);
\draw[gray,dashed,thin] (-1.3000000,-1.0000000)--(-1.3000000,1.0000000);
\node[above,fill=none,inner sep=2pt] at (-2.3000000,1.1500000) {角膜};
\node[above,fill=none,inner sep=2pt] at (-0.8000000,1.5000000) {晶状体};
\draw[gray,dashed,thin] (-0.8000000,1.4000000)--(-1.1000000,0.4500000);
\fill (2.4000000,0.0000000) circle (1.2pt);
\node[right,fill=none,inner sep=2pt] at (2.4000000,0.0000000) {视网膜};
\node[below,fill=none,inner sep=2pt] at (0.0000000,-1.5500000) {虚线平面表示眼的等效折射系统};
\end{tikzpicture}
\figcap{远物平行光在正常眼视网膜同一点会聚。光线按等效薄系统画，解剖轮廓用于辨认结构。}\label{ch28:fig:s43-20-3}
\end{minipage}\end{center}
}

\newcommand{\ApparentDepthDiagram}{%
\begin{center}\begin{minipage}{\linewidth}\centering
\begin{tikzpicture}[>=Latex,font=\small,every node/.style={fill=none}]
\begin{scope}[shift={(-3.0000000,0.0000000)}]
\draw[thick] (-2.0000000,0.0000000)--(2.0000000,0.0000000);
\fill (0.0000000,-1.8000000) circle (1.2pt);
\node[below,fill=none,inner sep=2pt] at (0.0000000,-1.8000000) {};
\draw[gray,dashed,thin] (0.0000000,-1.8000000)--(-0.7000000,-1.8000000);
\node[left,fill=none,inner sep=2pt] at (-0.7500000,-1.8000000) {$O$};
\draw[gray,dashed,thin] (0.0000000,-2.0000000)--(0.0000000,1.6000000);
\draw[DeepBlue,thick,->] (0.0000000,-1.8000000)--(-0.6000000,0.0000000);
\draw[DeepBlue,thick,->] (-0.6000000,0.0000000)--(-1.1902918,1.2694706);
\draw[gray,dashed,thin] (-0.6000000,0.0000000)--(0.0000000,-1.2903488);
\draw[DeepBlue,thick,->] (0.0000000,-1.8000000)--(0.6000000,0.0000000);
\draw[DeepBlue,thick,->] (0.6000000,0.0000000)--(1.1902918,1.2694706);
\draw[gray,dashed,thin] (0.6000000,0.0000000)--(0.0000000,-1.2903488);
\fill (0.0000000,-1.2903488) circle (1.2pt);
\node[below,fill=none,inner sep=2pt] at (0.0000000,-1.2903488) {};
\draw[gray,dashed,thin] (0.0000000,-1.2903488)--(0.7000000,-1.2903488);
\node[right,fill=none,inner sep=2pt] at (0.7500000,-1.2903488) {$I$};
\node[left,fill=none,inner sep=2pt] at (-1.3000000,-1.0000000) {水};
\node[left,fill=none,inner sep=2pt] at (-1.3000000,0.6000000) {空气中的眼};
\end{scope}
\begin{scope}[shift={(3.0000000,0.0000000)}]
\draw[thick] (-2.0000000,0.0000000)--(2.0000000,0.0000000);
\fill (0.0000000,-1.8000000) circle (1.2pt);
\node[below,fill=none,inner sep=2pt] at (0.0000000,-1.8000000) {};
\draw[gray,dashed,thin] (0.0000000,-1.8000000)--(-0.7000000,-1.8000000);
\node[left,fill=none,inner sep=2pt] at (-0.7500000,-1.8000000) {$O$};
\draw[gray,dashed,thin] (0.0000000,-2.0000000)--(0.0000000,1.6000000);
\draw[DeepBlue,thick,->] (0.0000000,-1.8000000)--(-0.6000000,0.0000000);
\draw[DeepBlue,thick,->] (-0.6000000,0.0000000)--(-0.9320392,1.3600551);
\draw[gray,dashed,thin] (-0.6000000,0.0000000)--(0.0000000,-2.4576411);
\draw[DeepBlue,thick,->] (0.0000000,-1.8000000)--(0.6000000,0.0000000);
\draw[DeepBlue,thick,->] (0.6000000,0.0000000)--(0.9320392,1.3600551);
\draw[gray,dashed,thin] (0.6000000,0.0000000)--(0.0000000,-2.4576411);
\fill (0.0000000,-2.4576411) circle (1.2pt);
\node[below,fill=none,inner sep=2pt] at (0.0000000,-2.4576411) {};
\draw[gray,dashed,thin] (0.0000000,-2.4576411)--(0.7000000,-2.4576411);
\node[right,fill=none,inner sep=2pt] at (0.7500000,-2.4576411) {$I$};
\node[left,fill=none,inner sep=2pt] at (-1.3000000,-1.0000000) {空气};
\node[left,fill=none,inner sep=2pt] at (-1.3000000,0.6000000) {水中的眼};
\end{scope}
\end{tikzpicture}
\figcap{从空气看水中（左）和从水中看空气（右）。虚像取所示两斜光线的反向交点；近轴极限为 $|v|=n_2u/n_1$。}\label{ch28:fig:s43-20-4}
\end{minipage}\end{center}
}

\newcommand{\ThinLensPanel}{%
\begin{center}\begin{minipage}{\linewidth}\centering
\begin{tikzpicture}[>=Latex,font=\small,every node/.style={fill=none}]
\begin{scope}[shift={(-5.0000000,0.0000000)}]
\draw[thick,fill=DeepBlue!5] (-.12,-1).. controls (-.65,-.3) and (-.65,.3)..(-.12,1)--(.12,1)..controls(.65,.3)and(.65,-.3)..(.12,-1)--cycle;
\node[below,fill=none,inner sep=2pt] at (0.0000000,-1.3000000) {双凸};
\end{scope}
\begin{scope}[shift={(-3.0000000,0.0000000)}]
\draw[thick,fill=DeepBlue!5] (0,-1)--(0,1)--(.12,1)..controls(.65,.3)and(.65,-.3)..(.12,-1)--cycle;
\node[below,fill=none,inner sep=2pt] at (0.0000000,-1.3000000) {平凸};
\end{scope}
\begin{scope}[shift={(-1.0000000,0.0000000)}]
\draw[thick,fill=DeepBlue!5] (-.1,-1)..controls(.3,-.3)and(.3,.3)..(-.1,1)--(.1,1)..controls(.9,.3)and(.9,-.3)..(.1,-1)--cycle;
\node[below,fill=none,inner sep=2pt] at (0.0000000,-1.3000000) {会聚弯月};
\end{scope}
\begin{scope}[shift={(1.0000000,0.0000000)}]
\draw[thick,fill=DeepBlue!5] (-.55,-1)..controls(-.05,-.3)and(-.05,.3)..(-.55,1)--(.55,1)..controls(.05,.3)and(.05,-.3)..(.55,-1)--cycle;
\node[below,fill=none,inner sep=2pt] at (0.0000000,-1.3000000) {双凹};
\end{scope}
\begin{scope}[shift={(3.0000000,0.0000000)}]
\draw[thick,fill=DeepBlue!5] (0,-1)--(0,1)--(.65,1)..controls(.1,.3)and(.1,-.3)..(.65,-1)--cycle;
\node[below,fill=none,inner sep=2pt] at (0.0000000,-1.3000000) {平凹};
\end{scope}
\begin{scope}[shift={(5.0000000,0.0000000)}]
\draw[thick,fill=DeepBlue!5] (-.45,-1)..controls(.45,-.3)and(.45,.3)..(-.45,1)--(.2,1)..controls(.6,.3)and(.6,-.3)..(.2,-1)--cycle;
\node[below,fill=none,inner sep=2pt] at (0.0000000,-1.3000000) {发散弯月};
\end{scope}
\begin{scope}[shift={(0.0000000,-4.0000000)}]
\draw[thick,fill=DeepBlue!5] (0.1669697,-0.8000000)--(0.1600696,-0.7840000)--(0.1533338,-0.7680000)--(0.1467607,-0.7520000)--(0.1403484,-0.7360000)--(0.1340954,-0.7200000)--(0.1280000,-0.7040000)--(0.1220607,-0.6880000)--(0.1162760,-0.6720000)--(0.1106446,-0.6560000)--(0.1051649,-0.6400000)--(0.0998358,-0.6240000)--(0.0946559,-0.6080000)--(0.0896241,-0.5920000)--(0.0847392,-0.5760000)--(0.0800000,-0.5600000)--(0.0754055,-0.5440000)--(0.0709546,-0.5280000)--(0.0666464,-0.5120000)--(0.0624799,-0.4960000)--(0.0584542,-0.4800000)--(0.0545684,-0.4640000)--(0.0508217,-0.4480000)--(0.0472133,-0.4320000)--(0.0437423,-0.4160000)--(0.0404082,-0.4000000)--(0.0372101,-0.3840000)--(0.0341475,-0.3680000)--(0.0312197,-0.3520000)--(0.0284260,-0.3360000)--(0.0257660,-0.3200000)--(0.0232390,-0.3040000)--(0.0208446,-0.2880000)--(0.0185823,-0.2720000)--(0.0164517,-0.2560000)--(0.0144522,-0.2400000)--(0.0125836,-0.2240000)--(0.0108454,-0.2080000)--(0.0092373,-0.1920000)--(0.0077591,-0.1760000)--(0.0064103,-0.1600000)--(0.0051907,-0.1440000)--(0.0041002,-0.1280000)--(0.0031385,-0.1120000)--(0.0023053,-0.0960000)--(0.0016006,-0.0800000)--(0.0010243,-0.0640000)--(0.0005761,-0.0480000)--(0.0002560,-0.0320000)--(0.0000640,-0.0160000)--(0.0000000,0.0000000)--(0.0000640,0.0160000)--(0.0002560,0.0320000)--(0.0005761,0.0480000)--(0.0010243,0.0640000)--(0.0016006,0.0800000)--(0.0023053,0.0960000)--(0.0031385,0.1120000)--(0.0041002,0.1280000)--(0.0051907,0.1440000)--(0.0064103,0.1600000)--(0.0077591,0.1760000)--(0.0092373,0.1920000)--(0.0108454,0.2080000)--(0.0125836,0.2240000)--(0.0144522,0.2400000)--(0.0164517,0.2560000)--(0.0185823,0.2720000)--(0.0208446,0.2880000)--(0.0232390,0.3040000)--(0.0257660,0.3200000)--(0.0284260,0.3360000)--(0.0312197,0.3520000)--(0.0341475,0.3680000)--(0.0372101,0.3840000)--(0.0404082,0.4000000)--(0.0437423,0.4160000)--(0.0472133,0.4320000)--(0.0508217,0.4480000)--(0.0545684,0.4640000)--(0.0584542,0.4800000)--(0.0624799,0.4960000)--(0.0666464,0.5120000)--(0.0709546,0.5280000)--(0.0754055,0.5440000)--(0.0800000,0.5600000)--(0.0847392,0.5760000)--(0.0896241,0.5920000)--(0.0946559,0.6080000)--(0.0998358,0.6240000)--(0.1051649,0.6400000)--(0.1106446,0.6560000)--(0.1162760,0.6720000)--(0.1220607,0.6880000)--(0.1280000,0.7040000)--(0.1340954,0.7200000)--(0.1403484,0.7360000)--(0.1467607,0.7520000)--(0.1533338,0.7680000)--(0.1600696,0.7840000)--(0.1669697,0.8000000)--(0.4330303,0.8000000)--(0.4399304,0.7840000)--(0.4466662,0.7680000)--(0.4532393,0.7520000)--(0.4596516,0.7360000)--(0.4659046,0.7200000)--(0.4720000,0.7040000)--(0.4779393,0.6880000)--(0.4837240,0.6720000)--(0.4893554,0.6560000)--(0.4948351,0.6400000)--(0.5001642,0.6240000)--(0.5053441,0.6080000)--(0.5103759,0.5920000)--(0.5152608,0.5760000)--(0.5200000,0.5600000)--(0.5245945,0.5440000)--(0.5290454,0.5280000)--(0.5333536,0.5120000)--(0.5375201,0.4960000)--(0.5415458,0.4800000)--(0.5454316,0.4640000)--(0.5491783,0.4480000)--(0.5527867,0.4320000)--(0.5562577,0.4160000)--(0.5595918,0.4000000)--(0.5627899,0.3840000)--(0.5658525,0.3680000)--(0.5687803,0.3520000)--(0.5715740,0.3360000)--(0.5742340,0.3200000)--(0.5767610,0.3040000)--(0.5791554,0.2880000)--(0.5814177,0.2720000)--(0.5835483,0.2560000)--(0.5855478,0.2400000)--(0.5874164,0.2240000)--(0.5891546,0.2080000)--(0.5907627,0.1920000)--(0.5922409,0.1760000)--(0.5935897,0.1600000)--(0.5948093,0.1440000)--(0.5958998,0.1280000)--(0.5968615,0.1120000)--(0.5976947,0.0960000)--(0.5983994,0.0800000)--(0.5989757,0.0640000)--(0.5994239,0.0480000)--(0.5997440,0.0320000)--(0.5999360,0.0160000)--(0.6000000,0.0000000)--(0.5999360,-0.0160000)--(0.5997440,-0.0320000)--(0.5994239,-0.0480000)--(0.5989757,-0.0640000)--(0.5983994,-0.0800000)--(0.5976947,-0.0960000)--(0.5968615,-0.1120000)--(0.5958998,-0.1280000)--(0.5948093,-0.1440000)--(0.5935897,-0.1600000)--(0.5922409,-0.1760000)--(0.5907627,-0.1920000)--(0.5891546,-0.2080000)--(0.5874164,-0.2240000)--(0.5855478,-0.2400000)--(0.5835483,-0.2560000)--(0.5814177,-0.2720000)--(0.5791554,-0.2880000)--(0.5767610,-0.3040000)--(0.5742340,-0.3200000)--(0.5715740,-0.3360000)--(0.5687803,-0.3520000)--(0.5658525,-0.3680000)--(0.5627899,-0.3840000)--(0.5595918,-0.4000000)--(0.5562577,-0.4160000)--(0.5527867,-0.4320000)--(0.5491783,-0.4480000)--(0.5454316,-0.4640000)--(0.5415458,-0.4800000)--(0.5375201,-0.4960000)--(0.5333536,-0.5120000)--(0.5290454,-0.5280000)--(0.5245945,-0.5440000)--(0.5200000,-0.5600000)--(0.5152608,-0.5760000)--(0.5103759,-0.5920000)--(0.5053441,-0.6080000)--(0.5001642,-0.6240000)--(0.4948351,-0.6400000)--(0.4893554,-0.6560000)--(0.4837240,-0.6720000)--(0.4779393,-0.6880000)--(0.4720000,-0.7040000)--(0.4659046,-0.7200000)--(0.4596516,-0.7360000)--(0.4532393,-0.7520000)--(0.4466662,-0.7680000)--(0.4399304,-0.7840000)--(0.4330303,-0.8000000)--(0.1669697,-0.8000000);
\draw[gray,thin,->] (-4.4000000,0.0000000)--(4.2666120,0.0000000);
\fill (-4.0000000,0.0000000) circle (1.2pt);
\node[above,fill=none,inner sep=2pt] at (-4.0000000,0.0000000) {$O$};
\fill (3.8166120,0.0000000) circle (1.2pt);
\node[above,fill=none,inner sep=2pt] at (3.8166120,0.0000000) {$I$};
\fill (0.0308631,0.3500000) circle (1.2pt);
\node[above left,fill=none,inner sep=2pt] at (0.0308631,0.3500000) {$P$};
\fill (0.5697641,0.3464526) circle (1.2pt);
\node[above right,fill=none,inner sep=2pt] at (0.5697641,0.3464526) {$Q$};
\draw[DeepBlue,thick,->] (-4.0000000,0.0000000)--(0.0308631,0.3500000);
\draw[DeepBlue,thick,->] (0.0308631,0.3500000)--(0.5697641,0.3464526);
\draw[DeepBlue,thick,->] (0.5697641,0.3464526)--(3.8166120,0.0000000);
\draw[gray,dashed,thin] (2.0000000,0.0000000)--(0.0308631,0.3500000);
\draw[gray,dashed,thin] (-1.4000000,0.0000000)--(0.5697641,0.3464526);
\fill (2.0000000,0.0000000) circle (1.2pt);
\node[below,fill=none,inner sep=2pt] at (2.0000000,0.0000000) {};
\fill (-1.4000000,0.0000000) circle (1.2pt);
\node[below,fill=none,inner sep=2pt] at (-1.4000000,0.0000000) {};
\draw[gray,dashed,thin] (2.0000000,0.0000000)--(2.0000000,-0.9000000);
\draw[gray,dashed,thin] (-1.4000000,0.0000000)--(-1.4000000,-0.9000000);
\node[below,fill=none,inner sep=2pt] at (2.0000000,-0.9500000) {$C_1$};
\node[below,fill=none,inner sep=2pt] at (-1.4000000,-0.9500000) {$C_2$};
\node[above,fill=none,inner sep=2pt] at (-1.2000000,1.0000000) {$C_1P=R_1$};
\node[above,fill=none,inner sep=2pt] at (2.2000000,1.0000000) {$C_2Q=|R_2|$};
\draw[gray,<->] (-4.0000000,-1.5000000)--(0.0000000,-1.5000000);
\node[below,fill=none,inner sep=2pt] at (-2.0000000,-1.5000000) {$u$};
\draw[gray,<->] (0.6000000,-1.5000000)--(3.8166120,-1.5000000);
\node[below,fill=none,inner sep=2pt] at (2.2083060,-1.5000000) {$v$};
\end{scope}
\end{tikzpicture}
\figcap{上为六种闭合透镜截面，置于空气中时左三种会聚、右三种发散；下为两球面分别折射的光路，$C_1P,C_2Q$ 是两面的半径和法线。}\label{ch28:fig:s43-20-5}
\end{minipage}\end{center}
}

\newcommand{\LensImagingDiagram}{%
\begin{center}\begin{minipage}{\linewidth}\centering
\begin{tikzpicture}[>=Latex,font=\small,every node/.style={fill=none}]
\draw[gray,thin,->] (-4.8000000,0.0000000)--(4.0000000,0.0000000);
\draw[DeepBlue,thick,<->] (0.0000000,-1.8500000)--(0.0000000,1.8500000);
\draw[black,thick,->] (-4.2000000,0.0000000)--(-4.2000000,1.5000000);
\node[above,fill=none,inner sep=2pt] at (-4.2000000,1.5000000) {物};
\draw[black,thick,->] (3.3123028,0.0000000)--(3.3123028,-1.1829653);
\node[below,fill=none,inner sep=2pt] at (3.3123028,-1.1829653) {像};
\fill (-3.7037037,0.0000000) circle (1.2pt);
\node[below,fill=none,inner sep=2pt] at (-3.7037037,0.0000000) {};
\node[below,fill=none,inner sep=2pt] at (-3.7037037,-0.5000000) {$2F$};
\fill (-1.8518519,0.0000000) circle (1.2pt);
\node[below,fill=none,inner sep=2pt] at (-1.8518519,0.0000000) {};
\node[below,fill=none,inner sep=2pt] at (-1.8518519,-0.5000000) {$F$};
\fill (1.8518519,0.0000000) circle (1.2pt);
\node[below,fill=none,inner sep=2pt] at (1.8518519,0.0000000) {};
\node[above,fill=none,inner sep=2pt] at (1.8518519,0.2000000) {$F'$};
\fill (3.7037037,0.0000000) circle (1.2pt);
\node[below,fill=none,inner sep=2pt] at (3.7037037,0.0000000) {};
\node[above,fill=none,inner sep=2pt] at (3.7037037,0.2000000) {$2F'$};
\draw[DeepBlue,thick,->] (-4.2000000,1.5000000)--(0.0000000,1.5000000);
\draw[DeepBlue,thick,->] (0.0000000,1.5000000)--(3.3123028,-1.1829653);
\draw[DeepBlue,thick,->] (-4.2000000,1.5000000)--(0.0000000,0.0000000);
\draw[DeepBlue,thick,->] (0.0000000,0.0000000)--(3.3123028,-1.1829653);
\draw[DeepBlue,thick,->] (-4.2000000,1.5000000)--(0.0000000,-1.1829653);
\draw[DeepBlue,thick,->] (0.0000000,-1.1829653)--(3.3123028,-1.1829653);
\end{tikzpicture}
\figcap{$u=70.0\,\mathrm{cm}$、$f\approx30.9\,\mathrm{cm}$、$v\approx55.2\,\mathrm{cm}$ 时的三条近轴特征光线；通过物方焦点的那条出射后平行光轴。}\label{ch28:fig:s43-20-6}
\end{minipage}\end{center}
}

\newcommand{\MagnifierDiagram}{%
\begin{center}\begin{minipage}{\linewidth}\centering
\begin{tikzpicture}[>=Latex,font=\small,every node/.style={fill=none}]
\begin{scope}[shift={(-2.6000000,0.0000000)}]
\draw[gray,thin,->] (-3.6000000,0.0000000)--(1.7000000,0.0000000);
\draw[DeepBlue,thick,<->] (0.0000000,-1.2000000)--(0.0000000,1.2000000);
\fill (-1.5000000,0.0000000) circle (1.2pt);
\node[below,fill=none,inner sep=2pt] at (-1.5000000,0.0000000) {$F$};
\draw[black,thick,->] (-1.0000000,0.0000000)--(-1.0000000,0.6500000);
\node[left,fill=none,inner sep=2pt] at (-1.0000000,0.6500000) {物};
\draw[black,thick,->] (-3.0000000,0.0000000)--(-3.0000000,1.9500000);
\node[above,fill=none,inner sep=2pt] at (-3.0000000,1.9500000) {虚像};
\draw[DeepBlue,thick,->] (-1.0000000,0.6500000)--(0.0000000,0.6500000);
\draw[DeepBlue,thick,->] (0.0000000,0.6500000)--(1.2000000,0.1300000);
\draw[gray,dashed,thin] (-3.0000000,1.9500000)--(0.0000000,0.6500000);
\draw[DeepBlue,thick,->] (-1.0000000,0.6500000)--(0.0000000,0.0000000);
\draw[DeepBlue,thick,->] (0.0000000,0.0000000)--(1.2000000,-0.7800000);
\draw[gray,dashed,thin] (-3.0000000,1.9500000)--(0.0000000,0.0000000);
\draw[Gold,very thick] (0.4000000,-0.3400000)--(0.4000000,0.5000000);
\node[right,fill=none,inner sep=2pt] at (0.5000000,1.2500000) {眼瞳位置};
\draw[gray,dashed,thin] (0.8000000,1.2000000)--(0.4000000,0.5000000);
\end{scope}
\begin{scope}[shift={(3.8000000,0.7000000)}]
\draw[black,thick,->] (-1.6000000,0.0000000)--(-1.6000000,0.6500000);
\draw[DeepBlue,thick,->] (-1.6000000,0.6500000)--(1.4000000,0.0000000);
\draw[gray,dashed,thin] (-1.6000000,0.0000000)--(1.4000000,0.0000000);
\node[right,fill=none,inner sep=2pt] at (1.4000000,0.0000000) {裸眼};
\draw[gray,<->] (-1.6000000,-0.4500000)--(1.4000000,-0.4500000);
\node[below,fill=none,inner sep=2pt] at (-0.1000000,-0.4500000) {$L_0$};
\node[above,fill=none,inner sep=2pt] at (0.0000000,1.2000000) {$\theta_0\simeq H/L_0$};
\end{scope}
\end{tikzpicture}
\figcap{左为放大镜的正立虚像及进入眼瞳的光；右为裸眼在明视距离观察同样大小物体的基准视角。}\label{ch28:fig:s43-20-7}
\end{minipage}\end{center}
}

\newcommand{\SphericalMirrorDiagram}{%
\begin{center}\begin{minipage}{\linewidth}\centering
\begin{tikzpicture}[>=Latex,font=\small,every node/.style={fill=none}]
\draw[thick] plot coordinates {(2.6488428,-1.4084147) (2.6624817,-1.3824585) (2.6758662,-1.3563702) (2.6889951,-1.3301523) (2.7018672,-1.3038074) (2.7144812,-1.2773379) (2.7268358,-1.2507464) (2.7389300,-1.2240354) (2.7507625,-1.1972075) (2.7623322,-1.1702652) (2.7736381,-1.1432111) (2.7846790,-1.1160479) (2.7954539,-1.0887780) (2.8059618,-1.0614041) (2.8162016,-1.0339288) (2.8261724,-1.0063547) (2.8358732,-0.9786845) (2.8453031,-0.9509208) (2.8544612,-0.9230663) (2.8633466,-0.8951236) (2.8719585,-0.8670954) (2.8802960,-0.8389843) (2.8883584,-0.8107932) (2.8961449,-0.7825245) (2.9036547,-0.7541811) (2.9108872,-0.7257657) (2.9178416,-0.6972809) (2.9245172,-0.6687295) (2.9309135,-0.6401143) (2.9370297,-0.6114379) (2.9428655,-0.5827031) (2.9484201,-0.5539126) (2.9536930,-0.5250692) (2.9586838,-0.4961756) (2.9633919,-0.4672347) (2.9678170,-0.4382491) (2.9719586,-0.4092216) (2.9758162,-0.3801551) (2.9793896,-0.3510522) (2.9826784,-0.3219158) (2.9856822,-0.2927487) (2.9884008,-0.2635536) (2.9908340,-0.2343333) (2.9929814,-0.2050906) (2.9948430,-0.1758284) (2.9964184,-0.1465493) (2.9977076,-0.1172563) (2.9987105,-0.0879520) (2.9994269,-0.0586393) (2.9998567,-0.0293211) (3.0000000,0.0000000) (2.9998567,0.0293211) (2.9994269,0.0586393) (2.9987105,0.0879520) (2.9977076,0.1172563) (2.9964184,0.1465493) (2.9948430,0.1758284) (2.9929814,0.2050906) (2.9908340,0.2343333) (2.9884008,0.2635536) (2.9856822,0.2927487) (2.9826784,0.3219158) (2.9793896,0.3510522) (2.9758162,0.3801551) (2.9719586,0.4092216) (2.9678170,0.4382491) (2.9633919,0.4672347) (2.9586838,0.4961756) (2.9536930,0.5250692) (2.9484201,0.5539126) (2.9428655,0.5827031) (2.9370297,0.6114379) (2.9309135,0.6401143) (2.9245172,0.6687295) (2.9178416,0.6972809) (2.9108872,0.7257657) (2.9036547,0.7541811) (2.8961449,0.7825245) (2.8883584,0.8107932) (2.8802960,0.8389843) (2.8719585,0.8670954) (2.8633466,0.8951236) (2.8544612,0.9230663) (2.8453031,0.9509208) (2.8358732,0.9786845) (2.8261724,1.0063547) (2.8162016,1.0339288) (2.8059618,1.0614041) (2.7954539,1.0887780) (2.7846790,1.1160479) (2.7736381,1.1432111) (2.7623322,1.1702652) (2.7507625,1.1972075) (2.7389300,1.2240354) (2.7268358,1.2507464) (2.7144812,1.2773379) (2.7018672,1.3038074) (2.6889951,1.3301523) (2.6758662,1.3563702) (2.6624817,1.3824585) (2.6488428,1.4084147)};
\draw[gray,thin,->] (-3.6000000,0.0000000)--(3.6000000,0.0000000);
\fill (-3.0000000,0.0000000) circle (1.2pt);
\node[below,fill=none,inner sep=2pt] at (-3.0000000,0.0000000) {$O$};
\fill (0.0000000,0.0000000) circle (1.2pt);
\node[below,fill=none,inner sep=2pt] at (0.0000000,0.0000000) {$C$};
\fill (3.0000000,0.0000000) circle (1.2pt);
\node[below right,fill=none,inner sep=2pt] at (3.0000000,0.0000000) {$V$};
\fill (2.9108872,0.7257657) circle (1.2pt);
\node[below,fill=none,inner sep=2pt] at (2.9108872,0.7257657) {};
\draw[gray,dashed,thin] (2.9108872,0.7257657)--(2.3608872,1.1757657);
\node[above left,fill=none,inner sep=2pt] at (2.3108872,1.2257657) {$P$};
\fill (1.0202029,0.0000000) circle (1.2pt);
\node[below,fill=none,inner sep=2pt] at (1.0202029,0.0000000) {$I$};
\draw[gray,dashed,thin] (0.0000000,0.0000000)--(2.9108872,0.7257657);
\draw[DeepBlue,thick,->] (-3.0000000,0.0000000)--(2.9108872,0.7257657);
\draw[DeepBlue,thick,->] (2.9108872,0.7257657)--(1.0202029,0.0000000);
\end{tikzpicture}
\figcap{球面镜真实反射光路，$CP$ 是法线。所示有限高度轴交点趋向近轴像点，但不与所有孔径光线严格重合。}\label{ch28:fig:s43-20-8}
\end{minipage}\end{center}
}

\newcommand{\AplanaticLensPanel}{%
% Independent panel: original physical paths and local coordinates retained.
\begin{center}\begin{minipage}{\linewidth}\centering
\begin{tikzpicture}[>=Latex,font=\small,every node/.style={fill=none}]
\draw[thick] (1.2122511,-0.9000000)--(1.2200945,-0.8775000)--(1.2277163,-0.8550000)--(1.2351183,-0.8325000)--(1.2423025,-0.8100000)--(1.2492707,-0.7875000)--(1.2560247,-0.7650000)--(1.2625661,-0.7425000)--(1.2688965,-0.7200000)--(1.2750174,-0.6975000)--(1.2809304,-0.6750000)--(1.2866368,-0.6525000)--(1.2921381,-0.6300000)--(1.2974353,-0.6075000)--(1.3025299,-0.5850000)--(1.3074230,-0.5625000)--(1.3121156,-0.5400000)--(1.3166089,-0.5175000)--(1.3209038,-0.4950000)--(1.3250014,-0.4725000)--(1.3289026,-0.4500000)--(1.3326081,-0.4275000)--(1.3361189,-0.4050000)--(1.3394356,-0.3825000)--(1.3425590,-0.3600000)--(1.3454899,-0.3375000)--(1.3482287,-0.3150000)--(1.3507762,-0.2925000)--(1.3531327,-0.2700000)--(1.3552990,-0.2475000)--(1.3572753,-0.2250000)--(1.3590622,-0.2025000)--(1.3606600,-0.1800000)--(1.3620690,-0.1575000)--(1.3632896,-0.1350000)--(1.3643219,-0.1125000)--(1.3651662,-0.0900000)--(1.3658228,-0.0675000)--(1.3662916,-0.0450000)--(1.3665729,-0.0225000)--(1.3666667,0.0000000)--(1.3665729,0.0225000)--(1.3662916,0.0450000)--(1.3658228,0.0675000)--(1.3651662,0.0900000)--(1.3643219,0.1125000)--(1.3632896,0.1350000)--(1.3620690,0.1575000)--(1.3606600,0.1800000)--(1.3590622,0.2025000)--(1.3572753,0.2250000)--(1.3552990,0.2475000)--(1.3531327,0.2700000)--(1.3507762,0.2925000)--(1.3482287,0.3150000)--(1.3454899,0.3375000)--(1.3425590,0.3600000)--(1.3394356,0.3825000)--(1.3361189,0.4050000)--(1.3326081,0.4275000)--(1.3289026,0.4500000)--(1.3250014,0.4725000)--(1.3209038,0.4950000)--(1.3166089,0.5175000)--(1.3121156,0.5400000)--(1.3074230,0.5625000)--(1.3025299,0.5850000)--(1.2974353,0.6075000)--(1.2921381,0.6300000)--(1.2866368,0.6525000)--(1.2809304,0.6750000)--(1.2750174,0.6975000)--(1.2688965,0.7200000)--(1.2625661,0.7425000)--(1.2560247,0.7650000)--(1.2492707,0.7875000)--(1.2423025,0.8100000)--(1.2351183,0.8325000)--(1.2277163,0.8550000)--(1.2200945,0.8775000)--(1.2122511,0.9000000)--(1.7860571,0.9000000)--(1.7972183,0.8775000)--(1.8080307,0.8550000)--(1.8185004,0.8325000)--(1.8286334,0.8100000)--(1.8384351,0.7875000)--(1.8479110,0.7650000)--(1.8570659,0.7425000)--(1.8659046,0.7200000)--(1.8744316,0.6975000)--(1.8826511,0.6750000)--(1.8905670,0.6525000)--(1.8981833,0.6300000)--(1.9055035,0.6075000)--(1.9125310,0.5850000)--(1.9192691,0.5625000)--(1.9257206,0.5400000)--(1.9318886,0.5175000)--(1.9377758,0.4950000)--(1.9433846,0.4725000)--(1.9487175,0.4500000)--(1.9537768,0.4275000)--(1.9585645,0.4050000)--(1.9630827,0.3825000)--(1.9673332,0.3600000)--(1.9713178,0.3375000)--(1.9750380,0.3150000)--(1.9784953,0.2925000)--(1.9816912,0.2700000)--(1.9846269,0.2475000)--(1.9873034,0.2250000)--(1.9897220,0.2025000)--(1.9918835,0.1800000)--(1.9937888,0.1575000)--(1.9954385,0.1350000)--(1.9968334,0.1125000)--(1.9979740,0.0900000)--(1.9988606,0.0675000)--(1.9994937,0.0450000)--(1.9998734,0.0225000)--(2.0000000,0.0000000)--(1.9998734,-0.0225000)--(1.9994937,-0.0450000)--(1.9988606,-0.0675000)--(1.9979740,-0.0900000)--(1.9968334,-0.1125000)--(1.9954385,-0.1350000)--(1.9937888,-0.1575000)--(1.9918835,-0.1800000)--(1.9897220,-0.2025000)--(1.9873034,-0.2250000)--(1.9846269,-0.2475000)--(1.9816912,-0.2700000)--(1.9784953,-0.2925000)--(1.9750380,-0.3150000)--(1.9713178,-0.3375000)--(1.9673332,-0.3600000)--(1.9630827,-0.3825000)--(1.9585645,-0.4050000)--(1.9537768,-0.4275000)--(1.9487175,-0.4500000)--(1.9433846,-0.4725000)--(1.9377758,-0.4950000)--(1.9318886,-0.5175000)--(1.9257206,-0.5400000)--(1.9192691,-0.5625000)--(1.9125310,-0.5850000)--(1.9055035,-0.6075000)--(1.8981833,-0.6300000)--(1.8905670,-0.6525000)--(1.8826511,-0.6750000)--(1.8744316,-0.6975000)--(1.8659046,-0.7200000)--(1.8570659,-0.7425000)--(1.8479110,-0.7650000)--(1.8384351,-0.7875000)--(1.8286334,-0.8100000)--(1.8185004,-0.8325000)--(1.8080307,-0.8550000)--(1.7972183,-0.8775000)--(1.7860571,-0.9000000)--(1.2122511,-0.9000000);
\draw[gray,thin,->] (-3.4000000,0.0000000)--(3.5000000,0.0000000);
\fill (-3.0000000,0.0000000) circle (1.2pt);
\node[below,fill=none,inner sep=2pt] at (-3.0000000,0.0000000) {$Q$};
\fill (-1.3333333,0.0000000) circle (1.2pt);
\draw[gray,thin](-1.3333333,.04)--(-1.3333333,.52);\node[above,fill=none,inner sep=2pt] at (-1.3333333,.56) {$C_1$};
\fill (0.0000000,0.0000000) circle (1.2pt);
\node[below,fill=none,inner sep=2pt] at (0.0000000,0.0000000) {};
\draw[gray,dashed,thin] (0.0000000,0.0000000)--(0.0000000,-0.6500000);
\node[below,fill=none,inner sep=2pt] at (0.0000000,-0.7000000) {$C_2$};
\draw[DeepBlue,thick,->] (-1.3333333,0.0000000)--(1.2948880,-0.6184277);
\draw[DeepBlue,thick,->] (1.2948880,-0.6184277)--(1.8540496,-0.7500000);
\draw[DeepBlue,thick,->] (1.8540496,-0.7500000)--(2.9411498,-0.9179680);
\draw[gray,dashed,thin] (-3.0000000,0.0000000)--(1.8540496,-0.7500000);
\draw[gray,dashed,thin] (0.0000000,0.0000000)--(1.8540496,-0.7500000);
\draw[DeepBlue,thick,->] (-1.3333333,0.0000000)--(1.2948880,0.6184277);
\draw[DeepBlue,thick,->] (1.2948880,0.6184277)--(1.8540496,0.7500000);
\draw[DeepBlue,thick,->] (1.8540496,0.7500000)--(2.9411498,0.9179680);
\draw[gray,dashed,thin] (-3.0000000,0.0000000)--(1.8540496,0.7500000);
\draw[gray,dashed,thin] (0.0000000,0.0000000)--(1.8540496,0.7500000);
\node[left,fill=none,inner sep=2pt] at (1.4000000,1.0500000) {$S_1$};
\node[right,fill=none,inner sep=2pt] at (2.1000000,1.0500000) {$S_2$};
\fill(1.8540496,.75)circle(1.2pt);\node[right]at(2.65,.4){$P$};\draw[thin](1.8540496,.75)--(2.55,.35);
\node[align=center]at(-.6,-1.6){$R_1$（左），$R_2$（右）\\玻璃折射率 $n$};
\end{tikzpicture}
\figcap{无球差透镜的 $Q$ 为虚像；虚线是出射光反向延长线，$C_1,C_2$ 是两球面曲率中心。}\label{ch28:fig:s43-20-9}
\end{minipage}\end{center}
}

\newcommand{\CatEyePanel}{%
% Independent panel: original physical paths and local coordinates retained.
\begin{center}\begin{minipage}{\linewidth}\centering
\begin{tikzpicture}[>=Latex,font=\small,every node/.style={fill=none}]
\draw[pattern=north east lines] (-1.8,.9) rectangle (1.8,1.6);\draw[pattern=north east lines] (-1.8,-1.6) rectangle (1.8,-.9);
\begin{scope}[shift={(-1.0000000,0.0000000)}]
\draw[thick,fill=DeepBlue!5] (-.55,-1)..controls(-.05,-.3)and(-.05,.3)..(-.55,1)--(.55,1)..controls(.05,.3)and(.05,-.3)..(.55,-1)--cycle;
\end{scope}
\begin{scope}[shift={(1.1000000,0.0000000)}]
\draw[thick,fill=DeepBlue!5] (-.12,-1).. controls (-.65,-.3) and (-.65,.3)..(-.12,1)--(.12,1)..controls(.65,.3)and(.65,-.3)..(.12,-1)--cycle;
\end{scope}
\node[below,fill=none,inner sep=2pt] at (-1.0000000,-2.5500000) {$f_1=-1\,\mathrm{cm}$};
\node[below,fill=none,inner sep=2pt] at (1.1000000,-2.5500000) {$f_2=3.5\,\mathrm{cm}$};
\node[left,fill=none,inner sep=2pt] at (-2.1000000,0.0000000) {室外};
\node[right,fill=none,inner sep=2pt] at (2.1000000,0.0000000) {室内};
\draw[gray,<->] (-1.0000000,-1.9500000)--(1.1000000,-1.9500000);
\node[below,fill=none,inner sep=2pt] at (0.0500000,-1.9500000) {$2.1\,\mathrm{cm}$};
\end{tikzpicture}
\figcap{猫眼的两透镜装在同一贯通门孔中；正向观察时光从室外凹透镜传向室内凸透镜。}\label{ch28:fig:s43-20-10}
\end{minipage}\end{center}
}

\newcommand{\SilveredLensPanel}{%
% Independent panel: original physical paths and local coordinates retained.
\begin{center}\begin{minipage}{\linewidth}\centering
\begin{tikzpicture}[>=Latex,font=\small,every node/.style={fill=none}]
\draw[gray,thin,->] (-3.5000000,0.0000000)--(0.5000000,0.0000000);
\draw[very thick,double] (0.0000000,-1.4000000)--(0.0000000,1.4000000);
\draw[black,thick,->] (-3.0000000,0.0000000)--(-3.0000000,0.9000000);
\draw[black,thick,->] (-1.0000000,0.0000000)--(-1.0000000,-0.3000000);
\draw[DeepBlue,thick,->] (-3.0000000,0.9000000)--(0.0000000,0.9000000);
\draw[DeepBlue,thick,->] (0.0000000,0.9000000)--(-1.0000000,-0.3000000);
\draw[DeepBlue,thick,->] (-3.0000000,0.9000000)--(0.0000000,0.0000000);
\draw[DeepBlue,thick,->] (0.0000000,0.0000000)--(-1.0000000,-0.3000000);
\node[above,fill=none,inner sep=2pt] at (-3.0000000,0.9000000) {$H$};
\node[below,fill=none,inner sep=2pt] at (-1.0000000,-0.5000000) {$-H/3$};
\node[above,fill=none,inner sep=2pt] at (0.0000000,1.5000000) {薄透镜与银面重合的等效平面};
\node[below,fill=none,inner sep=2pt] at (-2.0000000,-1.2000000) {$u=2f,\ v=2f/3$};
\end{tikzpicture}
\figcap{镀银薄透镜的等效折返光路，实像倒立，像高为物高的三分之一。}\label{ch28:fig:s43-20-11}
\end{minipage}\end{center}
}

\newcommand{\PrismTrainPanel}{%
% Independent panel: original physical paths and local coordinates retained.
\begin{center}\begin{minipage}{\linewidth}\centering
\begin{tikzpicture}[>=Latex,font=\small,every node/.style={fill=none}]
\draw[thick] (0.0000000,1.2000000)--(1.2000000,1.2000000)--(1.2000000,0.0000000)--(0.0000000,1.2000000);
\draw[DeepBlue,thick,->] (0.6000000,2.4000000)--(0.6000000,0.6000000);
\draw[DeepBlue,thick,->] (0.6000000,0.6000000)--(4.8000000,0.6000000);
\draw[black,thick,->] (0.6000000,2.4000000)--(0.9000000,2.4000000);
\node[above left,fill=none,inner sep=2pt] at (0.6000000,2.4000000) {$B$};
\node[above right,fill=none,inner sep=2pt] at (0.9000000,2.4000000) {$A$};
\begin{scope}[shift={(0.0000000,0.6000000)}]
\draw[DeepBlue,thick,<->] (2.4000000,-1.4000000)--(2.4000000,1.4000000);
\end{scope}
\draw[DeepBlue,thick] (4.2000000,-0.8000000)--(4.2000000,2.0000000);
\draw[DeepBlue,thick,->] (4.2000000,2.0000000)--(4.2000000,1.7000000);
\draw[DeepBlue,thick,->] (4.2000000,-0.8000000)--(4.2000000,-0.5000000);
\node[above,fill=none,inner sep=2pt] at (2.4000000,2.1500000) {$f_1=10$};
\node[above,fill=none,inner sep=2pt] at (4.2000000,2.1500000) {$f_2=-10$};
\draw[gray,<->] (-0.2500000,1.2000000)--(-0.2500000,2.4000000);
\node[left,fill=none,inner sep=2pt] at (-0.2500000,1.8000000) {$6$};
\draw[gray,<->] (-0.2500000,0.0000000)--(-0.2500000,1.2000000);
\node[left,fill=none,inner sep=2pt] at (-0.2500000,0.6000000) {$6$};
\draw[gray,<->] (1.2000000,-0.6000000)--(2.4000000,-0.6000000);
\node[below,fill=none,inner sep=2pt] at (1.8000000,-0.6000000) {$10$};
\draw[gray,<->] (2.4000000,-0.6000000)--(4.2000000,-0.6000000);
\node[below,fill=none,inner sep=2pt] at (3.3000000,-0.6000000) {$15$};
\node[below,fill=none,inner sep=2pt] at (2.4000000,-1.6000000) {长度单位：cm（横向比例压缩）};
\end{tikzpicture}
\figcap{累次成像的光轴在棱镜斜面上转过直角；物体 $AB$ 位于棱镜上方。}\label{ch28:fig:prism-train}
\end{minipage}\end{center}
}

\newcommand{\SequentialImagingDiagram}{%
\begin{center}\begin{minipage}{\linewidth}\centering
\begin{tikzpicture}[>=Latex,font=\small,every node/.style={fill=none}]
\draw[gray,thin,->] (-3.4000000,0.0000000)--(6.2000000,0.0000000);
\draw[DeepBlue,thick,<->] (0.0000000,-1.8000000)--(0.0000000,1.8000000);
\draw[DeepBlue,thick,<->] (4.0000000,-1.8000000)--(4.0000000,1.8000000);
\draw[black,thick,->] (-3.0000000,0.0000000)--(-3.0000000,0.6000000);
\node[above,fill=none,inner sep=2pt] at (-3.0000000,0.6000000) {物 $O$};
\draw[black,thick,->] (1.5000000,0.0000000)--(1.5000000,-0.3000000);
\draw[gray,dashed,thin] (1.5000000,-0.3000000)--(1.5000000,1.0000000);
\node[above,fill=none,inner sep=2pt] at (1.5000000,1.0500000) {中间像 $I_1$};
\draw[black,thick,->] (5.6666667,0.0000000)--(5.6666667,0.2000000);
\node[above,fill=none,inner sep=2pt] at (5.6666667,0.2000000) {最终像 $I_2$};
\draw[DeepBlue,thick,->] (-3.0000000,0.6000000)--(0.0000000,0.4000000);
\draw[DeepBlue,thick,->] (0.0000000,0.4000000)--(4.0000000,-1.4666667);
\draw[DeepBlue,thick,->] (4.0000000,-1.4666667)--(5.6666667,0.2000000);
\draw[DeepBlue,thick,->] (-3.0000000,0.6000000)--(0.0000000,0.0000000);
\draw[DeepBlue,thick,->] (0.0000000,0.0000000)--(4.0000000,-0.8000000);
\draw[DeepBlue,thick,->] (4.0000000,-0.8000000)--(5.6666667,0.2000000);
\node[above,fill=none,inner sep=2pt] at (0.0000000,1.9000000) {$f_1$};
\node[above,fill=none,inner sep=2pt] at (4.0000000,1.9000000) {$f_2$};
\draw[gray,<->] (0.0000000,-2.0000000)--(4.0000000,-2.0000000);
\node[below,fill=none,inner sep=2pt] at (2.0000000,-2.0000000) {$D$};
\end{tikzpicture}
\figcap{两透镜逐级成像：$I_1$ 同时是第二透镜的物，$u_2=D-v_1$；实际光线穿过 $I_1$ 时没有折转。\label{ch28:fig:sequential}}
\end{minipage}\end{center}
}

% Source revision 2026-09-23: physical optics; 12 cm print width, >=9 pt.
% Layout: two independent horizontal return paths, f=1.5, L=3.
% [S at front focus -> lens -> mirror] | [S -> lens -> image on mirror].
% Rays are blue; return rays gold; dashed lines are axes only. Transparent labels.
\newcommand{\SelfConjugateDiagram}{%
\begin{center}\begin{minipage}{\linewidth}\centering\begin{tikzpicture}[>=Latex,font=\small,scale=.88]
\begin{scope}
 \draw[gray,dashed](-1.8,0)--(3.3,0);
 \draw[thick,<->](0,-1.2)--(0,1.2);\draw[very thick](3,-1.2)--(3,1.2);
 \fill(-1.5,0)circle(1.4pt);\node[above left=4pt]at(-1.5,0){$S=F$};
 \foreach \h in {-.65,.65}{
  \draw[DeepBlue,thick,->](-1.5,0)--(0,\h)--(3,\h);
  \draw[Gold!80!black,thick,->](2.6,\h)--(.2,\h);
 }
 \node[above]at(0,1.3){透镜};\node[above]at(3,1.3){平面镜};
 \draw[<->](-1.5,-1.5)--(0,-1.5)node[midway,below]{$f$};
 \node at(.65,-2.2){(a) 平行光自准直};
\end{scope}
\begin{scope}[xshift=7cm]
 \draw[gray,dashed](-3.3,0)--(3.3,0);
 \draw[thick,<->](0,-1.2)--(0,1.2);\draw[very thick](3,-1.2)--(3,1.2);
 \fill(-3,0)circle(1.4pt)node[below]{$S$};\fill(3,0)circle(1.4pt)node[right]{$I_1$};
 \draw[DeepBlue,thick,->](-3,0)--(0,.65)--(3,0);
 \draw[Gold!80!black,thick,->](3,0)--(0,-.65)--(-3,0);
 \draw[<->](-3,-1.5)--(0,-1.5)node[midway,below]{$u$};
 \draw[<->](0,-1.5)--(3,-1.5)node[midway,below]{$L$};
 \node[above]at(0,1.3){透镜};\node[above]at(3,1.3){平面镜};
 \node at(0,-2.2){(b) 第一次像落在镜面};
\end{scope}
\end{tikzpicture}\figcap{透镜和平面镜的两种自共轭。右图用 $L=2f$ 示范，返程光线与另一条去程光线互换。\label{ch28:fig:self-conjugate}}\end{minipage}\end{center}}


几乎所有光学系统都有一个共同任务：收集从物点发出的部分光线，经过反射或折射后，使它们严格或近似会聚于另一个点；或者使出射光看起来像从某一点发出。从波动观点说，光学元件改变入射波面的形状，使各部分在像点具有相同相位。完美成像一般需要复杂非球面和极高加工精度；实际系统往往限制孔径和视场，并在像质、通光量、体积与成本之间折中。

这项任务包含两个不可混淆的要求：一是“同一个物点”的各条光线应在同一个像点重合，二是不同物点的像还应保持可识别的相对位置。前者若不满足会形成模糊斑，后者若不满足会产生畸变。增加孔径能收集更多光、提高衍射分辨率，却会让球面像差更显著；缩小孔径可改善近轴近似，却降低亮度。所谓光学系统的像质，正是在这些条件和加工误差共同作用下评定，而不是只看某一条特征光线是否穿过焦点。

球面容易加工的原因也很直接：两个相配曲面在研磨时，高出的区域优先被磨去，长期相对运动会趋向共同球面。非球面能够更好校正像差，但面形检测和装调更难。图\ref{ch28:fig:contexts}A 的研磨情景解释为什么实际器件偏爱球面；后面的图~\ref{ch28:fig:s43-20-1} 则说明使用球面所付出的代价：一般只能在一定孔径内近似成像。

实际设计还必须限定视场。轴上物点的近轴光也许能很好会聚，离轴物点却可能出现彗差、像散和像面弯曲；因此“一个焦距”只描述系统的一阶性质，并不等于整幅画面处处清晰。镜头组通过不同材料和不同曲率的元件互相补偿这些误差，代价是表面数量、重量、装调难度和成本增加。

成像结论必须同时注明孔径、视场和近轴条件。

\historyfigure{optical-instruments-built-in-v1.png}{成像器件与仪器：光具座、暗箱、透镜、显微镜、望远镜和人眼。}

图\ref{ch28:fig:contexts}把加工、折射现象、成像实验和仪器放在同一条学习线上。A 的玻璃坯与研磨盘以配合球面接触；B 的刚性鸭形模型跨越水面，水下部分的视位置改变；C 的光源、透镜和屏共用一条光具座；D 的反光镜将镜头来光折向上方取景器。下面分别用球面、平面界面和逐级成像解释这些现象，场景图不用于量取角度或焦距。
\historyfigure[0.94\linewidth]{imaging-contexts.png}{研磨球面（A）、水面折射演示（B）、透镜成像实验（C）和单反相机剖面（D）。\label{ch28:fig:contexts}}



\subsection{成像不是画一条“代表光线”}
一点发出的光通常占据许多方向。只有这些光线经过系统后在另一点附近会聚，或其反向延长线在另一点附近相交，才可以谈这个物点的像。只追踪一条光线不能确定像的位置；至少要取两条独立光线，并检查邻近光线是否也满足同一成像关系。

所谓近轴，是相对于每个元件自身主轴而言，入射高度和倾角都足够小，使 $\sin\theta\simeq\tan\theta\simeq\theta$。它不等于“整张图画得很小”。一条相对图纸水平线倾角很小的光线，在强倾斜的曲面上仍可能大角度入射。使用球面或薄透镜公式时，既要检查光线相对于轴的角度，也要检查孔径相对于曲率半径是否足够小。

本章的物距与像距采用实物、实像为正的约定。折射后出射光真实相交处为实像；出射光发散而其反向延长线相交处为虚像。前一元件的像若位于后一元件的下游，后一元件截获的是尚未会聚的光束，这个点对后一元件而言就是虚物。先判断会聚、发散和传播方向，再赋予符号，比根据“左负右正”机械代数更可靠。


\section{平面界面与视深}



平面界面相当于 $R\to\infty$ 的球面折射面，故
\begin{equation}
 \frac{n_1}{u}+\frac{n_2}{v}=0,\qquad v=-\frac{n_2}{n_1}u.
\end{equation}
实物在单一平面折射后成为虚像。观察者在低折射率介质中看高折射率介质里的物体，视深为真实深度乘 $n_2/n_1$，所以水底看起来变浅；反过来从水中看空气里的物体，视距离增大。结论只对近轴光线成立，大视角时不同光线的反向延长线不严格共点。

图~\ref{ch28:fig:s43-20-4} 可直接给出这一比例。物体在介质 $n_1$ 中的真实深度为 $u$，观察者在介质 $n_2$ 中。取一条靠近法线的光，横向位移 $y$ 同时满足
\[
 y\simeq u\theta_1\simeq |v|\theta_2,
 \qquad n_1\theta_1\simeq n_2\theta_2.
\]
消去角度得
\[
 |v|=\frac{n_2}{n_1}u.
\]
从空气看水中，$n_2/n_1\approx1/1.33$，所以视深约为真实深度的四分之三；从水中看空气，比例反过来。这个推导同时说明观察方向若偏离法线很大，$\tan\theta\simeq\theta$ 与 $\sin\theta\simeq\theta$ 不再成立，所谓唯一“视深”只是近似。

\subsection{为何视深公式只能近轴使用}
设物点在平界面下深度 $h$ 处，物方折射率为 $n_1$，观察侧折射率为 $n_2$。选一条到达界面的光线，与法线夹角分别为 $i,r$，其界面交点离轴距离为 $\rho=h\tan i$。把出射光向后延长到轴，它与轴的交点深度为
\[
 h'=\frac{\rho}{\tan r}=h\frac{\tan i}{\tan r}.
\]
结合 $n_1\sin i=n_2\sin r$，有
\[
 \frac{h'}h=\frac{n_2\cos r}{n_1\cos i}
 \simeq\frac{n_2}{n_1}\quad(i,r\ll1).
\]
在水下向空气看，$n_1>n_2$，所以近轴视深小于真实深度。大角度时余弦比不可略去，不同方向的光线延长后与轴的交点不同，因而不能说平界面在任意孔径下都形成一个严格的点像。

对多层平行透明介质，逐层相加的是横向位移。若第 $j$ 层厚度为 $t_j$、折射率为 $n_j$，最终观察介质的折射率为 $n_o$，近轴时 $n_j\theta_j=n_o\theta_o$，故
\[
 \rho=\sum_jt_j\theta_j
 =\theta_o\sum_j\frac{n_ot_j}{n_j},
 \qquad h_{\rm app}=\sum_j\frac{n_ot_j}{n_j}.
\]
观察者到最外界面的距离若也计入总视距离，须另加这一段；它不是某一层介质的视厚度。全部 $n_j=n_o$ 时，视距离应回到几何距离，这是检验层数、起点和终点有没有算错的直接办法。


\ApparentDepthDiagram



\Needspace{8\baselineskip}\begin{origexamplebox}{水中直杆的最大表观弯折}
\label{ch28:ex:depth-stick}
一根直杆与水平液面成角 $\alpha$，液体对空气的折射率为 $n>1$。从近于正上方观察，采用近轴视深模型，杆在液体内的像与杆的真实方向相差 $\beta$。求使 $\beta$ 最大的 $\alpha$，并求最大弯折角。
\end{origexamplebox}\par\nopagebreak[3]
\begin{xesolution}
这里求的是观察到的杆像方向，不是某一条折射光线转过的角度。图~\ref{ch28:fig:depth-stick} 中杆与液面交于 $O$，杆上一点 $P$ 的水平距离为 $x$、真实深度为 $h$。正上方的近轴观察使横向位置保持不变，而深度变为 $h'=h/n$，所以它的虚像在 $P'$。

% Figure design: geometric optics; original drawing, no reference-image trace.
% Print width <=11cm; 9-10pt labels, transparent. One vertical section.
% Nodes O=(0,0), Q=(2.2,0), P=(2.2,-3.11127), P'=(2.2,-1.55563).
% Example geometry n=2, tan(alpha)=sqrt(2); h/h'=2 exactly.
% Solid OP=real rod, dashed OP'=apparent rod; vertical dimensions on right.
% Angles: alpha at r=.65, alpha' at r=1.20, beta between rays at r=1.8.
% Minimal figure, no decorative panels. Whitespace above water for view arrow.
% ASCII: observer up | O----Q----water; O\P' (dashed); O\P (solid); h,h' right.
\begin{center}\begin{minipage}{\linewidth}\centering
\begin{tikzpicture}[>=Latex,every node/.style={fill=none,font=\small,inner sep=2pt}]
 \coordinate (O) at (0,0);\coordinate (Q) at (2.2,0);
 \coordinate (P) at (2.2,{-2.2*sqrt(2)});
 \coordinate (V) at (2.2,{-1.1*sqrt(2)});
 \draw[thick] (-.6,0)--(4.45,0);
 \node[left] at (-.6,0){液面};
 \draw[very thick,DeepBlue] (O)--(P);
 \draw[thick,dashed,DeepBlue] (O)--(V);
 \draw[guide] (Q)--(P);
 \fill (O) circle(1.3pt);\fill (P) circle(1.3pt);\fill (V) circle(1.3pt);
 \node[above left] at (O){$O$};\node[above left] at (Q){$Q$};
 \node[right=5pt] at (P){$P$};\node[right=5pt] at (V){$P'$};
 \draw[->] (2.2,.4)--(2.2,1.05);
 \node[above] at (2.2,1.12){近于正上方观察};
 \draw (.65,0) arc[start angle=0,end angle=-54.73561,radius=.65];
 \node at (-18:.91){$\alpha$};
 \draw (1.2,0) arc[start angle=0,end angle=-35.26439,radius=1.2];
 \node at (-13:1.56){$\alpha'$};
 \draw (-35.26439:1.8) arc[start angle=-35.26439,end angle=-54.73561,radius=1.8];
 \node at (-45:2.12){$\beta$};
 \draw[<->] (3.2,0)--(3.2,{-2.2*sqrt(2)});
 \node[right] at (3.2,-2.45){$h$};
 \draw[<->] (4.05,0)--(4.05,{-1.1*sqrt(2)});
 \node[right] at (4.05,-.78){$h'$};
 \node[below] at (1.1,-3.65){实线为杆，虚线为杆的像，不是光线};
\end{tikzpicture}
\figcap{虚像 $P'$ 比 $P$ 浅，水平坐标不变。图取 $n=2$，角度关系由同一组点确定。\label{ch28:fig:depth-stick}}
\end{minipage}\end{center}

记像杆与液面的角为 $\alpha'$，则
\[
 \tan\alpha=\frac hx,\qquad
 \tan\alpha'=\frac{h'}x=\frac{\tan\alpha}{n},\qquad
 \beta=\alpha-\arctan\frac{\tan\alpha}{n}.
\]
对 $0<\alpha<\pi/2$ 求导时使用弧度。令 $t=\tan\alpha$，有
\[
 \frac{\dd\beta}{\dd\alpha}
 =1-\frac{n(1+t^2)}{n^2+t^2}
 =\frac{(n-1)(n-t^2)}{n^2+t^2}.
\]
分母为正且 $n>1$，故 $t^2<n$ 时递增，$t^2>n$ 时递减。唯一最大值发生在
\[
 \boxed{\alpha=\arctan\sqrt n},\qquad
 \alpha'=\arctan\frac1{\sqrt n}=\frac\pi2-\alpha.
\]
由此
\[
 \boxed{\beta_{\max}=2\arctan\sqrt n-\frac\pi2
 =\arcsin\frac{n-1}{n+1}}.
\]
当杆趋于水平或竖直时，弯折角都趋于零；当 $n\to1$ 时也趋于零。以水的 $n=4/3$ 为例，$\alpha\approx49.11^\circ$、$\beta_{\max}\approx8.21^\circ$。结论建立在近于正上方、局部近轴的观察模型上，不能用来描述任意斜视时整根杆的严格像。
\end{xesolution}


\section{球面折射与无球差点}



透镜是最常见的成像元件，人眼本身也包含折射系统。严格点到点成像通常要求非球面，但球面易于研磨加工；在限制孔径、只取近轴光线后，球面可以近似成像。

图~\ref{ch28:fig:s43-20-1} 中，物点 $S$ 和像点 $P$ 位于球面顶点 $V$ 的两侧，顶点距离分别为 $u,v$；球面半径为 $R$，两侧折射率为 $n_1,n_2$。若入射点 $A$ 对应球心角 $\varphi$，两段几何长度为
\[
 \ell_1=\sqrt{(u+R)^2+R^2-2R(u+R)\cos\varphi},
\]
\[
 \ell_2=\sqrt{(v-R)^2+R^2+2R(v-R)\cos\varphi}.
\]
上述几何先针对 $u,v,R>0$ 的实物实像配置：令 $V=(0,0)$、$C=(R,0)$，则 $A=(R-R\cos\varphi,R\sin\varphi)$，$S=(-u,0)$、$P=(v,0)$。直接计算 $|A-S|$ 和 $|P-A|$ 就得到两式；特别地，$\varphi=0$ 时分别回到 $u,v$。
总光程 $\mathcal L=n_1\ell_1+n_2\ell_2$。对非轴上入射点求导，费马条件给出
\[
 \frac{\partial\mathcal L}{\partial\varphi}
 =R\sin\varphi\left[\frac{n_1(u+R)}{\ell_1}-\frac{n_2(v-R)}{\ell_2}\right]=0.
\]
括号内一般随 $\varphi$ 变化，故不能用一个固定的 $v$ 同时满足所有孔径光线。近轴条件 $|\varphi|\ll1$ 下，令 $\ell_1\to u,\ell_2\to v$，得到
$n_1(1+R/u)=n_2(1-R/v)$，即
\begin{equation}
 \frac{n_1}{u}+\frac{n_2}{v}=\frac{n_2-n_1}{R}.
\end{equation}
将近轴式推广到虚物和虚像时，采用“实正虚负”的物像距约定，而不是把左右坐标都按同一方向读数：发散入射光对应实物 $u>0$，会聚入射光对应虚物 $u<0$；出射光真实会聚对应实像 $v>0$，出射光发散而反向延长线相交对应虚像 $v<0$。$R$ 在球心位于出射侧时为正，位于入射侧时为负。虚共轭须使用有向光程或延长线，不能把上面的正几何长度直接当成负长度。只要在每个界面重新识别物方、像方和传播方向，这一约定可连续使用。

近轴式并不是把严格关系中的角度直接“删掉”，而是保留光程随孔径高度 $y$ 的最低非平凡项。球面在顶点附近的弧高为 $x\simeq y^2/(2R)$。物点到界面的距离与界面到像点的距离分别展开为
\[
 \ell_1\simeq u+x+\frac{y^2}{2u},\qquad
 \ell_2\simeq v-x+\frac{y^2}{2v}.
\]
代入 $\mathcal L=n_1\ell_1+n_2\ell_2$ 后，若所有近轴光线到达像点的光程相同，$y^2$ 项的系数必须为零：
\[
 \frac{n_1}{2u}+\frac{n_2}{2v}+\frac{n_1-n_2}{2R}=0,
\]
这就得到球面折射公式。被略去的 $y^4,y^6,\ldots$ 项正是大孔径时球差等像差的来源，所以“球面能够成像”必须连同近轴、窄孔径条件一起表述。

严格关系中，$\ell_1,\ell_2$ 本身随入射点角度 $\varphi$ 改变；因此即使对某一条光线找到了令光程驻定的入射点，也不能保证所有 $\varphi$ 都给出同一个 $P$。近轴展开的作用，是让最低阶光程差只由 $u,v,R,n_1,n_2$ 决定而不再显含孔径高度。若要让大角度光也严格会聚，就必须改变曲面形状，使从物点经界面到像点的总光程对整个有效口径都相等；这一般给出非球面。例~\ref{ch28:ex:s43-20-6} 的特殊共轭点则利用几何与折射率关系，在球面上实现了少见的无球差成像。

\SphereRefractionPanel
\VirtualSignPanel

虚像点没有真实能量会聚，不能在该点直接放胶片或 CCD 接收，只能让另一光学系统或眼睛再次处理这些发散光。虚物则不会由普通发光物体单独产生，它只会出现在多元件系统中：前一元件产生的会聚光尚未到达会聚点，就先遇到下一界面。

“虚”描述的是到达当前元件之前或离开当前元件之后的光束状态，而不是说这个点没有物理意义。实物发出的是发散束；虚物对应尚未完成会聚的入射束。实像处确有光能会聚，可以放屏接收；虚像处没有出射光通过，只能由反向延长线定位。判断顺序应当是先看箭头方向，再看光束是会聚还是发散，最后决定 $u,v$ 的正负。仅凭某一点画在元件左边还是右边，在反射系统或折返光路中很容易出错。

例如，把一个新的折射面插入原本会在其后方会聚的光束中，这个会聚点就是该折射面的虚物，$u<0$；折射后若光仍在后方真实相交，则 $v>0$。反过来，普通物体经放大镜后出射光发散，延长线在物方相交，所以 $u>0,v<0$。图~\ref{ch28:fig:s43-20-2} 左图为实物经平界面形成虚像，右图为虚物经平界面形成实像。虚线分别表示出射光的反向延长线和未遇界面时的入射光会聚路径，不是真实传播光线。

\Needspace{8\baselineskip}\begin{origexamplebox}{【球透镜】}\label{ch28:ex:s43-20-1}
空气中半径为 $R$、折射率为 $n>1$ 的玻璃球，在近轴、窄孔径条件下把球外的小明成一个球外实像。求小明到其像的最小距离。
\end{origexamplebox}\par\nopagebreak[3]
\begin{xesolution}
设物到第一球面的距离为 $x$。第一次折射（空气到玻璃，半径 $R$）给出
\[
 \frac1x+\frac n{v_1}=\frac{n-1}{R},\qquad
 v_1=\frac{Rn x}{(n-1)x-R}.
\]
第二球面的物距为 $u_2=2R-v_1$，其有向半径为 $-R$。再次折射后得到球外实像距第二面
\[
 v_2=-R\frac{(n-2)x-2R}{2(n-1)x+(n-2)R}.
\]

%<optics-review:ball-second-surface>
第二次折射的代入过程可以展开为
\[
 \frac n{u_2}+\frac1{v_2}=\frac{1-n}{-R}=\frac{n-1}{R},
 \qquad
 u_2=R\frac{(n-2)x-2R}{(n-1)x-R}.
\]
于是
\[
 \frac1{v_2}=\frac{n-1}{R}
 -\frac{n[(n-1)x-R]}{R[(n-2)x-2R]}
 =-\frac{2(n-1)x+(n-2)R}{R[(n-2)x-2R]},
\]
取倒数即得上述 $v_2$。如果 $v_1>2R$，第一面折射后的光尚未会聚就遇到第二面，故第二面的物为虚物、$u_2<0$；不能把 $u_2$ 改成绝对值。

%</optics-review:ball-second-surface>
小明到像的总距离化简为
\[
 L=x+2R+v_2
 =\frac{2(n-1)(R+x)^2}{(n-2)R+2(n-1)x}.
\]
记分母 $D_x=(n-2)R+2(n-1)x$。实物及球外实像要求 $x>0$、$v_2>0$：当 $1<n<2$ 时，$x>(2-n)R/[2(n-1)]$；当 $n=2$ 时任意 $x>0$；当 $n>2$ 时，$0<x<2R/(n-2)$。
这些区间内均有 $D_x>0$。逐项求导得
\[
 L'(x)=\frac{4(n-1)(R+x)\{D_x-(n-1)(R+x)\}}{D_x^2}
 =\frac{4(n-1)(R+x)[(n-1)x-R]}{D_x^2}.
\]
所以 $x<R/(n-1)$ 时递减，$x>R/(n-1)$ 时递增，而驻点在上述每个允许区间内。因此全局最小值为
\[
 \boxed{L_{\min}=\frac{2nR}{n-1}}.
\]
此时光路关于球心对称；两次折射的符号均已按各自传播方向重新定义。

%<optics-review:ball-parallel-limit>
还需检查一个容易漏掉的特殊情况：最优物距 $x_*=R/(n-1)$ 恰使 $v_1$ 表达式的分母为零。这不是物理过程失效，而是第一面出射光变为平行光，中间像在无穷远。直接令第一面公式中的 $n/v_1=0$；第二面也有 $n/u_2=0$，因而
\[
 \frac1{v_2}=\frac{n-1}{R},\qquad
 v_2=x_*=\frac R{n-1}.
\]
所以 $L=x_*+2R+v_2=2nR/(n-1)$，正好补上中间代数除法没有覆盖的这一点。$n\to1^+$ 时最小距离趋于无穷大，符合玻璃球与空气折射率相同时不再具有聚光能力；整个结果仍限于近轴、有效孔径较小的成像模型。

%</optics-review:ball-parallel-limit>
\end{xesolution}

% Relocated AplanaticLensPanel to its own solution.

\Needspace{8\baselineskip}\begin{origexamplebox}{【大角度成像】}\label{ch28:ex:s43-20-6}
显微镜物镜中有一枚两面均为球面的弯月形厚透镜。左表面 $S_1$ 球心为 $C_1$、半径 $R_1$；右表面 $S_2$ 球心为 $C_2$、半径 $R_2$。两球心都在相应表面的左侧，$C_1$ 又在 $C_2$ 左侧；玻璃对空气折射率为 $n>1$，且 $C_1C_2=R_2/n$。物位于 $C_1$。证明：（1）所有通过有效口径的出射光线的反向延长线严格交于一点 $Q$；（2）$QC_2=nR_2$，并判断像的虚实。
\end{origexamplebox}\par\nopagebreak[3]
\begin{xesolution}
\AplanaticLensPanel

从 $C_1$ 发出的任意光线到达 $S_1$ 时沿该球面的半径入射，所以第一面入射角为零，不发生偏折。于是玻璃内所有光线仍像从 $C_1$ 发出。

把图~\ref{ch28:fig:s43-20-9} 中第二球心取为原点：$C_2=(0,0)$、$C_1=(-R_2/n,0)$，在同侧选 $Q=(-nR_2,0)$。设第二面点 $P=R_2(\cos\varphi,\sin\varphi)$，并记 $a=|P-C_1|$、$b=|P-Q|$。直接计算
\[
 a^2=R_2^2(1+n^{-2}+2n^{-1}\cos\varphi),\qquad
 b^2=R_2^2(1+n^2+2n\cos\varphi)=n^2a^2.
\]
因此 $b=na$。外法向为 $\boldsymbol N=P/R_2$，入射方向 $(P-C_1)/a$，候选出射方向 $(P-Q)/b$。用二维叉积求它们与法线夹角的正弦：
\[
 \sin i=\frac{R_2|\sin\varphi|}{na},\qquad
 \sin r=\frac{nR_2|\sin\varphi|}{b}
 =\frac{R_2|\sin\varphi|}{a}=n\sin i.
\]
两方向的切向分量同号，符合玻璃到空气的折射定律。取通常朝右的有效口径 $|\varphi|<\pi/2$ 时，两方向与外法线的点积均为正，确为向外传播且不发生全反射；还需选取两球面间有正厚度、入射线先穿过 $S_1$ 的口径。
出射方向是从 $Q$ 指向 $P$ 再向右延伸，\textbf{不是从 $P$ 射向 $Q$}。故只有反向延长线在 $Q$ 相交，形成严格无球差虚像，且
\[
 \boxed{QC_2=nR_2}.
\]
证明没有使用小角近似，所以 $C_1,Q$ 是该球面的两个共轭无球差点；后续物镜可利用这束仿佛从 $Q$ 发出的光继续成像。图中到 $Q$ 的虚线表示延长线，在 $Q$ 处不能直接放屏接收此像。
\end{xesolution}

\section{球面镜、薄透镜与物像关系}



非球面反射镜造价较高，球面反射镜在近轴范围内同样可近似成像。图~\ref{ch28:fig:s43-20-8} 中 $CP$ 是法线，反射角等于入射角；对三角形作小角近似可得
\begin{equation}
 \frac1u+\frac1v=\frac2R=\frac1f,\qquad f=\frac R2.
\end{equation}
实物、虚物、实像和虚像的符号意义与透镜一致，只需注意反射后传播方向反转。横向放大率仍为 $m=-v/u$。凸球面镜可取 $R<0$、$f<0$，其作用类似发散透镜。

几何推导中，入射点 $P$、球心 $C$ 和像点必须位于同一子午平面。设物点、像点和球心在主轴上的距离分别为 $u,v,R$，而 $P$ 相对轴的高度很小。半径 $CP$ 是法线，反射定律给出法线两侧角度相等。把三角形中的正弦用小角本身代替，并用“物点方向角 $+$ 像点方向角 $=2\times$ 法线方向角”，就得到
\[
 \frac yu+\frac yv\simeq2\frac yR.
\]
约去非零的 $y$ 后即为球面镜公式。这个过程明确显示了因子 $2$ 来自法线把方向转过两倍，而不是来自光程“往返两次”。

平面镜是 $R\to\infty$ 的极限，因此
\begin{equation}
 v=-u.
\end{equation}
普通实物经平面镜成等大的虚像；一束原本将会聚于镜后虚物点的光，经平面镜也可以在镜前形成实像。单镜头反光相机的折转光路中就包含类似过程。

图\ref{ch28:fig:contexts}D 中反光镜处于镜头预定焦面之前，它把仍会聚的光束折向取景光路，并不是把镜头的会聚能力消除；拍摄时反光镜让开，光才能到达后方感光面。

球面镜公式同样只在近轴范围成立。对凹面镜，平行轴的边缘光与近轴光通常不在同一点会聚，这是球差；把孔径缩小能减轻球差，却同时降低通光量并增强衍射。凹面镜的 $f=R/2$ 只是顶点附近的傍轴焦距，不应把图~\ref{ch28:fig:s43-20-8} 上任意高度的反射点都强行连到这一点。平面镜则没有这一曲率近似：用镜像法可精确证明物、像关于镜面对称，但虚像一侧并无真实光线通过。

镜像法的精确性可由任意一点验证。把物点 $O$ 关于镜面对称到 $O'$，对镜面上任一反射点 $P$ 都有 $PO=PO'$；反射定律又使折线 $O\to P\to$ 观察者等价于从 $O'$ 发出的直线。所有反射光的反向延长线因此严格通过同一个 $O'$，且 $|v|=u$、横向倍率为 $1$。若入射光原本向镜后一点会聚，把该点视为虚物并作同样的对称，反射光会在镜前真实会聚，这说明平面镜也能“虚物成实像”。

\SphericalMirrorDiagram



传统透镜至少包含一个曲面。中心比边缘厚的透镜通常会聚，中心较薄的透镜通常发散。忽略厚度的薄透镜可看作两个球面折射面的组合。对置于空气中的折射率 $n$ 的薄透镜，两个有向曲率半径为 $R_1,R_2$，连续应用球面折射式并消去中间像，得到制镜者公式和成像公式
\begin{equation}
 \frac1f=(n-1)\left(\frac1{R_1}-\frac1{R_2}\right),
 \qquad
 \frac1u+\frac1v=\frac1f.
\end{equation}
会聚透镜 $f>0$，发散透镜 $f<0$。横向放大率（带方向）为 $m=-v/u$，其绝对值
\begin{equation}
 M_T=\left|\frac vu\right|=\left|\frac f{u-f}\right|.
\end{equation}
负号表示实像倒立。作图时可用三条近轴特征光线：过光心近似不偏折；入射平行轴的光经像方焦点；入射通过物方焦点的光出射平行轴。

把两个折射面的中间步骤写出，可以避免把 $R_2$ 的符号背错。第一面由空气进入玻璃：
\[
 \frac1u+\frac n{v_1}=\frac{n-1}{R_1}.
\]
薄透镜厚度忽略后，第一面的像是第二面的虚物，故第二面物距 $u_2=-v_1$。从玻璃射回空气：
\[
 \frac n{u_2}+\frac1v=\frac{1-n}{R_2}.
\]
两式相加时 $v_1$ 恰好消去，才得到制镜者公式。双凸透镜按从左向右传播的约定通常有 $R_1>0,R_2<0$，所以两项共同贡献正光焦度；若只把两个半径都当作无符号长度，第二项就会误消掉第一项。

透镜成像图还提供一套数值自检：$u>2f$ 时为缩小倒立实像；$u=2f$ 时等大；$f<u<2f$ 时为放大倒立实像；$0<u<f$ 时为正立放大虚像。$u\to\infty$ 应给出 $v\to f$，$u\to f^+$ 时 $v\to+\infty$。这些极限若与计算结果不符，应先检查物距、焦距和传播方向的符号，而不是修改作图去迎合错误答案。

“中心厚”与“会聚”是对普通空气中、普通正折射率玻璃的经验说法；决定光焦度的仍是折射率差和两个有向曲率。把同一玻璃透镜浸入折射率更高的液体，光焦度可能减小、为零甚至反号。薄透镜近似还要求厚度相对曲率半径和物像距离足够小，使两个主平面可合并为一个平面；厚透镜中必须保留两面的间隔，中间像位置也不能简单取 $u_2=-v_1$。

图~\ref{ch28:fig:s43-20-5} 下图把两个界面分开：$P$ 处以 $C_1P$ 为法线折射，玻璃内部的 $PQ$ 保持直线；到 $Q$ 才以 $C_2Q$ 为法线再次折射。上面的薄透镜公式，是把这两个依次发生的过程在厚度可忽略时合并，而不是让光线在玻璃内部任意拐弯。
\ThinLensPanel
\LensImagingDiagram
图~\ref{ch28:fig:s43-20-6} 的三条特征光线在同一像点会合；通过物方焦点的光线出射后平行于光轴，不能把它误画成继续通过像方焦点。

图\ref{ch28:fig:contexts}C 的实验中，移动屏直到光斑最清晰，所测的屏距应与物距一起满足成像公式；若像为虚像，不能靠移动出射侧屏找到它。

例如玻璃折射率 $n=1.54$，双凸透镜两面半径大小分别为 $100\,\mathrm{cm}$ 和 $20\,\mathrm{cm}$（按方向取 $R_1=100\,\mathrm{cm},R_2=-20\,\mathrm{cm}$），则 $f\approx30.9\,\mathrm{cm}$。物距 $u=70.0\,\mathrm{cm}$ 时，$v=fu/(u-f)\approx55.2\,\mathrm{cm}$，$|m|\approx0.788$，形成缩小倒立的实像。

横向放大率来自过光心近似不偏折的两相似三角形，故带符号形式为 $m=-v/u$。它只描述垂直主轴的微小线段，不保证任意平面图形经透镜后保持相似：图形上不同点若有不同轴向坐标，就具有不同物距和放大率。圆、抛物线以及运动轨迹经透镜后是否仍保持原曲线类型，需要把每个点分别按 $X=fx/(x-f)$、$Y=-fy/(x-f)$ 变换，不能把一个统一倍率套到整幅有深度的图形上。这正是例~\ref{ch28:ex:s43-20-3} 的出发点。



\Needspace{8\baselineskip}\begin{origexamplebox}{【透镜成像】}\label{ch28:ex:s43-20-3}
半径为 $R$ 的圆经焦距为 $f$ 的凸透镜成像。（1）圆心放在哪里时像仍为圆？（2）抛物线能否仍成抛物线？（3）透镜中看到的抛体运动能否仍像抛体运动？
\end{origexamplebox}\par\nopagebreak[3]
\begin{xesolution}
取光轴为 $x$ 轴，物点 $(x,y)$ 的近轴像为
\[
 X=\frac{fx}{x-f},\qquad Y=-\frac{fy}{x-f}.
\]
这里先讨论图形位于包含光轴的子午面、圆心在轴上的情形。$x>0$ 表示物点位于透镜左侧距离 $x$ 处，$X>0$ 表示像点在右侧；两侧横向坐标均向上为正。反解得 $x=fX/(X-f)$、$y=-fY/(X-f)$。圆 $(x-a)^2+y^2=R^2$ 变换后为
\[
 [(f-a)X+af]^2+f^2Y^2=R^2(X-f)^2.
\]
展开后 $X^2$ 的系数为 $(f-a)^2-R^2$，$Y^2$ 的系数为 $f^2$，两者相等且非零才仍为圆，得到
\[
 (a-f)^2=R^2+f^2.
\]
对透镜同侧的实物圆，取
\[
 \boxed{a=f+\sqrt{f^2+R^2}};
\]
另一代数根描述另一侧或虚物配置。

对轴向抛物线 $y^2=4p(x-a)$，同样代入得
\[
 f^2Y^2=4p[(f-a)X+af](X-f).
\]
当 $a=f$ 时右侧化为 $4pf^2(X-f)$，所以 $Y^2=4p(X-f)$：\textbf{存在特殊放置}使抛物线仍成抛物线。一般放置则会变成椭圆或双曲型圆锥曲线。这里讨论各曲线上满足近轴与有效口径的有限部分，不意味着无穷大的图形都能通过透镜。

抛体的\textbf{轨迹形状}与\textbf{运动规律}要分开。若子午面内 $x(t)=x_0+v_xt$，则
\[
 X(t)=f+\frac{f^2}{x_0+v_xt-f},\qquad
 X''(t)=\frac{2f^2v_x^2}{(x_0+v_xt-f)^3}.
\]
当 $v_x\ne0$ 时，纵向像加速度随时间变化，即使像轨迹恰为抛物线，也一般不是匀加速度抛体。若圆心偏离光轴为 $(a,b)$，代入后的圆锥曲线另含 $2bfXY$ 项；要在这个子午面仍为圆，还须 $b=0$。这补足了第一问中圆心的横向位置条件。

题目若不限定图形朝向，还存在另一类简单答案：将圆或抛体轨迹放在\textbf{垂直光轴的同一物平面}内，所有点物距均为 $u_0$，倍率 $m=-f/(u_0-f)$ 为常数。圆自然仍成圆，无须前述特殊轴向位置。若该平面中的抛体坐标为
\[
 y(t)=y_0+v_yt,\qquad z(t)=z_0+v_zt-\tfrac12gt^2,
\]
像坐标为 $Y=my,Z=mz$，故 $Y''=0,Z''=-mg$。忽略光传播时间差，在近轴有效视场内，像确实仍作非退化的匀加速度抛体运动。因此第三问“能否”的回答是\textbf{能}，上述横向平面就是明确例子；其等效加速度通常不等于当地重力加速度。
\end{xesolution}

\subsection{从平行光到近物：七种入射光束的像}
式 $1/u+1/v=1/f$ 不只给数值，也控制出射光束是会聚、平行还是发散。图~\ref{ch28:fig:lens-states-real}、图~\ref{ch28:fig:lens-states-virtual} 将同一会聚薄透镜的七种状态逐一画出。光沿左向右传播，$F,F'$ 分别是物方、像方焦点；图中物点和像点都在轴上。蓝色实线是实际光线，灰色虚线仅是出射光的反向延长线，不能把虚线看成有光沿它反向传播。

% Figure contract: seven axial bundles, deterministic paraxial lens, f=1.2.
% Purpose: distinguish real/virtual/infinite images, restore seven source states.
% Print width 14cm; small text (9-10pt), no opaque label backgrounds.
% Shared coordinates: lens x=0, incident slope y/u, outgoing slope y/u-y/f.
% Layout: four real-image rows, then three parallel/divergent rows, 2.35cm pitch.
% Diagram x=-4.8..4.8; annotations x=5.0; ticks and labels below beam region.
% Rays are straight inside each medium; only the lens changes direction.
% ASCII: object * --> | lens | --> * image ; row title right, focal ticks below.
% For u=0 sentinel use parallel incidence; u=f is calculated without 1/(u-f).
% Quality: common f and aperture; seven cases are not decorative variants.
\newcommand{\LensStateRow}[5]{%
\begin{scope}[yshift=#1cm]
 \draw[gray,thin](-4.8,0)--(4.7,0);
 \draw[thick,<->](0,-.55)--(0,.55);
 \foreach \x/\lab in {-2.4/2F,-1.2/F,1.2/F',2.4/2F'}{
  \draw[gray] (\x,-.04)--(\x,.04);
  \draw[gray,thin] (\x,-.6)--(\x,-.73);
  \node[below,font=\scriptsize]at(\x,-.73){$\lab$};
 }
 \foreach \rawheight in {-.32,-.16,.16,.32}{
  \pgfmathsetmacro{\h}{ifthenelse(and(#2>0,#2<.6),.6875*\rawheight,\rawheight)}
  \pgfmathsetmacro{\slopein}{ifthenelse(#2==0,0,\h/max(#2,.001))}
  \pgfmathsetmacro{\slopeout}{\slopein-\h/1.2}
  \pgfmathsetmacro{\leftx}{ifthenelse(#2==0,-4.6,-#2)}
  \draw[DeepBlue,thin,->](\leftx,{\h+\slopein*\leftx})--(0,\h);
  \draw[DeepBlue,thin,->](0,\h)--(#3,{\h+\slopeout*#3});
  \ifnum#4=1
   \pgfmathsetmacro{\virtualx}{-\h/\slopeout}
   \draw[gray,dashed,thin](\virtualx,0)--(0,\h);
  \fi
 }
 \node[right,align=left,text width=3.6cm]at(5,.03){#5};
\end{scope}}
\begin{center}\begin{minipage}{\linewidth}\centering
\begin{tikzpicture}[>=Latex,every node/.style={fill=none,inner sep=2pt,font=\small},scale=.9]
 \LensStateRow{0}{0}{2.1}{0}{(a) $u\to\infty$\\平行入射，$v=f$}
 \LensStateRow{-2.35}{3.6}{3.2}{0}{(b) $u>2f$\\$f<v<2f$}
 \LensStateRow{-4.7}{2.4}{3.6}{0}{(c) $u=2f$\\$v=2f$}
 \LensStateRow{-7.05}{1.8}{4.6}{0}{(d) $f<u<2f$\\$v>2f$}
\end{tikzpicture}
\figcap{四种形成实像的轴上光束。所有会聚点都由 $v=fu/(u-f)$ 决定，而不是凭图形外观选取。\label{ch28:fig:lens-states-real}}
\end{minipage}\end{center}

图~\ref{ch28:fig:lens-states-real}(b)—(d) 中，物点靠近焦点时，像点向右远离透镜。代入 $u=3f,2f,3f/2$，分别得 $v=3f/2,2f,3f$。当 $u\to f^+$，$v\to+\infty$；因此不能在有限屏幕上继续寻找一个焦点附近的像。

\begin{center}\begin{minipage}{\linewidth}\centering
\begin{tikzpicture}[>=Latex,every node/.style={fill=none,inner sep=2pt,font=\small},scale=.9]
 \LensStateRow{0}{1.2}{4.6}{0}{(e) $u=f$\\出射平行\\像在无穷远}
 \LensStateRow{-2.35}{.9}{2.5}{1}{(f) $f/2<u<f$\\$v<0,\ |v|>f$}
 \LensStateRow{-4.7}{.55}{1.8}{1}{(g) $0<u<f/2$\\$v<0,\ |v|<f$}
\end{tikzpicture}
\figcap{从焦点移向透镜，出射光由平行变为发散；后两图的虚像由反向延长线相交确定。\label{ch28:fig:lens-states-virtual}}
\end{minipage}\end{center}

图~\ref{ch28:fig:lens-states-virtual}(f)、(g) 分开了虚像在物方焦点之外和之内的两种情况。由
\[
 |v|=\frac{fu}{f-u},\qquad
 |v|>f\ \Longleftrightarrow\ u>\frac f2
\]
可见分界不是第二种透镜，而是同一透镜的物距 $u=f/2$；这时虚像恰在 $F$。例如 $u=3f/4$ 时 $v=-3f$，$u=f/3$ 时 $v=-f/2$。虚像仍位于物点的左侧，但光并未真的会聚在那里。实际有限高度物体的正倒与大小，需再用横向放大率 $m=-v/u$ 判断：(b)—(d) 的实像为倒立，(f)、(g) 的虚像为正立放大。

\section{组合系统与逐次成像}



多元件系统应沿传播方向逐级成像：前一元件的像就是后一元件的物。两薄透镜焦距为 $f_1,f_2$、间距为 $D$，第一透镜物距为 $u_1$ 时
\[
 v_1=\frac{f_1u_1}{u_1-f_1},\qquad
 u_2=D-v_1,\qquad
 v_2=\frac{f_2u_2}{u_2-f_2}.
\]
$v_1>D$ 时，第二透镜前的光正在会聚，故 $u_2<0$ 是虚物；$v_1<0$ 时第一透镜产生虚像，第二透镜看到的物距为 $D+|v_1|$。总横向放大率是各级带符号放大率之积。只背数值而不画传播方向，很容易把虚物和虚像混淆。

\SequentialImagingDiagram

图\ref{ch28:fig:sequential}的 $I_1$ 位于两镜之间，所以第二透镜的物距为正；若把第二透镜移到 $I_1$ 左侧，$I_1$ 就成为第二透镜的虚物。这个正负转换来自光束状态，而不是把中间像删掉。

逐级计算时建议在每个元件前写三行：光线此刻向哪边传播、入射束是发散还是会聚、若没有该元件它将从哪里来或到哪里去。第一问决定正方向，第二问决定实物或虚物，第三问给出物距的几何大小。中间像即使落在两个元件之外，也不能被删除；它仍通过入射束的会聚或发散状态影响下一次成像。最后还应检查总放大率 $m=m_1m_2\cdots$ 的符号与图中倒立次数是否一致。

显微镜和望远镜体现了两种典型组合方式。显微镜的正物镜把近处小物形成放大的中间实像，正目镜再把这个中间像当作焦内实物，给眼睛形成放大的虚像；望远镜的物镜把远处目标形成焦面附近的中间实像，目镜主要放大它对眼睛的视角。两者都需要中间像进入目镜的有效孔径，否则即使成像公式给出位置，观察者也收不到完整光束。伽利略望远镜用负目镜在物镜光束尚未会聚前截取虚物，系统短且成正像，却没有可放屏或十字丝的中间实像，视场也较受限制。

评价组合系统不能只看倍率。物镜口径决定收集光的能力和衍射分辨率，焦距组合决定角放大率，孔径光阑决定哪些倾斜光束能通过，而像差、镀膜和装调决定这些理论性能能否实现。倍率超过物镜能够提供的空间信息时，只是把衍射斑和像差一同放大。于是“高级望远镜”的简易判断应先看有效口径、像质和机械稳定性，再看可用倍率，而不能只比较广告中的最大倍数。

\Needspace{14\baselineskip}





\subsection{中间像、虚物和无穷远像的衔接}
两个薄透镜沿同一方向依次排列，间距为 $d$。第一镜求出的像距为 $v_1$ 时，第二镜的物距是 $u_2=d-v_1$，不是 $d+v_1$。当 $v_1>d$，第一镜单独形成的像位于第二镜之后，入射到第二镜的光尚在会聚，于是 $u_2<0$，恰好对应虚物。

这个关系来自同一个轴向坐标：第一镜在 $x=0$，第二镜在 $x=d$，中间像在 $x=v_1$。因此第二镜到它的有向物距为 $d-v_1$。每一次代入前都把这一位置关系写出，能避免“中间像是实像，所以对下一镜也必是实物”的错误。

若物恰在第一镜物方焦平面上，第一镜出射光平行，不能把 $v_1=\infty$ 当普通有限数计算。轴外物点高度为 $h$ 时，出射光束倾角为 $\theta=-h/f_1$。第二镜把这束平行光聚到自身后焦平面，高度
\[
 h'=f_2\theta=-\frac{f_2}{f_1}h.
\]
这里通过“位置—角度—位置”完成成像，虽然中间像在无穷远，横向倍率仍是有限的 $-f_2/f_1$。两镜间距影响有限孔径能否接收整个光束，但不改变这个理想焦平面到焦平面的近轴倍率。

反射折返系统还多一步：反射之后要沿新的传播方向重新确定下一个物距。推荐先沿实际光路给元件编号，再逐次记录“当前元件—入射束会聚点—出射束会聚点”，最后才相乘各级横向倍率。屏的位置、像的位置、透镜的位置是三种不同对象，不能用同一个字母代替。

\Needspace{8\baselineskip}\begin{origexamplebox}{【奇妙的成像】}\label{ch28:ex:s43-20-2}
用焦距为 $f$ 的凸透镜和平面镜组合成像。点光源经透镜、镜面反射并再次通过透镜后，像回到光源。光源位于主轴上，平面镜垂直于主轴；若平面镜到透镜的距离为 $L$，求光源到透镜的两个可能距离，并说明真实光源何时可实现。
\end{origexamplebox}\par\nopagebreak[3]
\begin{xesolution}
把平面镜后的返程光路关于镜面展开，问题等价于两枚相同薄透镜相距 $2L$，物和最终像分别位于系统两端、距各自透镜均为 $u$。用空气中光线列向量 $(y,\theta)^{\mathsf T}$，传播矩阵为 $T(d)=\left(\begin{smallmatrix}1&d\\0&1\end{smallmatrix}\right)$，薄透镜矩阵为 $F=\left(\begin{smallmatrix}1&0\\-1/f&1\end{smallmatrix}\right)$。依次相乘得
\[
 FT(2L)F=\begin{pmatrix}1-2L/f&2L\\2(L-f)/f^2&1-2L/f\end{pmatrix}.
\]
再在左右各乘 $T(u)$，总矩阵的右上元为
\[
 B=2u(1-2L/f)+2L+\frac{2(L-f)u^2}{f^2}.
\]

%<optics-review:matrix-meaning>
为什么检查右上元 $B$ 就够了？从轴上物点发出一条倾角为 $\theta_0$ 的光，其初始列向量为 $(0,\theta_0)^{\mathsf T}$。若总矩阵为 $\left(\begin{smallmatrix}A&B\\C&D\end{smallmatrix}\right)$，到候选像面时光线高度便是 $y=B\theta_0$。不同 $\theta_0$ 都必须给出 $y=0$，所以必须 $B=0$。例如将上面的两镜矩阵记为 $\left(\begin{smallmatrix}a&b\\c&a\end{smallmatrix}\right)$，左右各乘传播矩阵后，右上元就是 $b+2au+cu^2$；代入 $a=1-2L/f$、$b=2L$、$c=2(L-f)/f^2$，便得到刚才的二次式。这里的每一项都有长度量纲。

%</optics-review:matrix-meaning>
同一物点发出的不同倾角光必须在同一像点重合，所以 $B=0$。乘去分母并分解，回到原点的条件即
\[
 (u-f)\,[Lf-Lu+fu]=0.
\]
故
\[
 \boxed{u_1=f,\qquad u_2=\frac{Lf}{L-f}}.
\]
第一种情况下透镜把轴上点源的光变为平行于轴的光，垂直放置的平面镜令其原路返回。第二种可不用矩阵理解：第一次透镜成像恰在镜面上，即 $v_1=L$，于是 $1/u+1/L=1/f$。镜面反射后，这个镜面上的像点成为返程实物，再经透镜成像回原点。两条不同的光线在镜面相交后互换返程路径，不能说每一条都沿自身原路返回。
\SelfConjugateDiagram
图~\ref{ch28:fig:self-conjugate} 中 $S=F$ 表示物点恰在透镜的物方焦点；第一次出射为平行光束，折返后再由透镜会聚。
第二根对应有限共轭位置。$u_2$ 要成为镜前有限距离实物位置需 $L>f$；$L<f$ 时第二根对应虚物，不能靠移动真实点光源实现；$L=f$ 时第二根趋于无穷远，有限物距方程只剩 $u=f$。
\end{xesolution}

\Needspace{8\baselineskip}\begin{origexamplebox}{【猫眼探秘】}\label{ch28:ex:s43-20-7}
门上的猫眼由室外侧焦距 $-1.0\,\mathrm{cm}$ 的凹透镜和室内侧焦距 $3.5\,\mathrm{cm}$ 的凸透镜组成，间距 $2.1\,\mathrm{cm}$。说明为什么室外看不清室内而室内能看清室外；若试图从室外看室内应怎样做？
\end{origexamplebox}\par\nopagebreak[3]
\begin{xesolution}
\CatEyePanel

%<optics-review:cat-eye-reference>
图~\ref{ch28:fig:s43-20-10} 标出了同一门孔中的室外凹透镜与室内凸透镜。以下先以远处景物作近轴计算；正向光按“凹透镜—凸透镜”传播，反向则按“凸透镜—凹透镜”传播。每到一个元件都沿新的传播方向判定实物与虚物，不能只交换两个焦距而保留原来的物距符号。

%</optics-review:cat-eye-reference>

室外远物先经凹透镜，成为位于凹透镜室外侧 $1.0\,\mathrm{cm}$ 的虚像。对凸透镜，这个物的距离为 $u=2.1+1.0=3.1\,\mathrm{cm}<3.5\,\mathrm{cm}$，故
\[
 v=\frac{3.5\times3.1}{3.1-3.5}\approx-27.1\,\mathrm{cm}.
\]
凸透镜把前一级的小虚像再放大，最终虚像位于可调焦的约 $27.1\,\mathrm{cm}$ 处。这里的“放大”针对中间像，并不表示门外景物的总角放大率大于一；广角门镜通常压缩景物的角尺度，实际视场还受孔径和像差限制。

反向观察远处室内景物时，凸透镜先把平行光会聚于其室外侧 $3.5\,\mathrm{cm}$ 处；凹透镜位于途中，面对的是其外侧 $1.4\,\mathrm{cm}$ 处的虚物，即 $u=-1.4\,\mathrm{cm}$。由 $f=-1.0\,\mathrm{cm}$ 得 $v=-3.5\,\mathrm{cm}$，反向光束在凹透镜后发散，反向延长线交于室内侧距该凹透镜 $3.5\,\mathrm{cm}$ 处。能观察到多大范围，还须检查门孔、两片透镜和眼瞳对斜光束的截取，不能仅凭像距断言视场大小。

反向观察还有更直接的调焦困难：若眼睛几乎贴着室外透镜，虚像到眼睛只有约 $3.5\,\mathrm{cm}$，远小于通常约 $25\,\mathrm{cm}$ 的近点；仅仅贴近孔径、遮挡杂散光不能让正常眼清楚调焦。眼距为 $e$ 时虚像距眼约 $e+3.5\,\mathrm{cm}$，后退可能满足近点条件，却会严重损失视场和接收光量。
从光学原理上，可用适当会聚光学系统重新成像，使最终像处于观察眼可调焦的范围；例如忽略元件间距时，让焦距约 $3.5\,\mathrm{cm}$ 的会聚元件以该虚像为实物，可使出射光近似平行。这只解决调焦，不消除门镜原有的有限孔径、像差与视场限制。由此可见，\textbf{瞳孔接光}和\textbf{眼睛调焦}是两个必须分别满足的条件。
\end{xesolution}

% Relocated SilveredLensPanel to its own solution.

\Needspace{8\baselineskip}\begin{origexamplebox}{【镀银透镜】}\label{ch28:ex:s43-20-8}
一平凸透镜焦距为 $f$，其平面镀银。凸面一侧距透镜 $2f$ 处放置高为 $H$ 的物，物垂直主轴且下端在轴上。（1）作图判断像的虚实；（2）计算像的位置和大小。
\end{origexamplebox}\par\nopagebreak[3]
\begin{xesolution}
\SilveredLensPanel

光先通过透镜，遇平面镀银层反射，再反向通过同一透镜。薄元件近似下，透镜光焦度 $1/f$ 被使用两次，平面镜本身不增加光焦度，所以组合等效为焦距
\[
 F=\frac f2
\]
的凹面镜。图~\ref{ch28:fig:s43-20-11}用等效薄平面画两条代表光线，它们返回物方并真实相交，故为实像。逐级追踪也得到同一结论：第一次透镜给 $v_1=2f$；在光尚未会聚前由紧贴的平面镜反射，会聚点关于镜面对称到物方 $2f$ 处。这个点对返程透镜是虚物，$u_2=-2f$，于是 $-1/(2f)+1/v=1/f$，得 $v=2f/3$。第一次倍率 $-1$、平面镜横向倍率 $+1$、返程透镜倍率 $+1/3$，乘积为 $-1/3$。

等效镜的物距 $u=2f=4F$，成像公式给出
\[
 \frac1{2f}+\frac1v=\frac1F=\frac2f,\qquad
 \boxed{v=\frac{2f}{3}}.
\]
像位于凸面外物方 $2f/3$ 处。横向放大率 $m=-v/u=-1/3$，所以
\[
 \boxed{H'=\frac H3},
\]
方向倒立。把镀银面误当成“成一次平面镜虚像”而停止追踪，会漏掉返程的第二次折射。
\end{xesolution}

\Needspace{22\baselineskip}\begin{origexamplebox}{【累次成像】}\label{ch28:ex:s43-20-9}
折射率 $n=1.5$、直角边长 $6\,\mathrm{cm}$ 的全反射棱镜上方 $6\,\mathrm{cm}$ 放置物体 $AB$。棱镜右侧 $10\,\mathrm{cm}$ 处为焦距 $f_1=10\,\mathrm{cm}$ 的凸透镜，再向右 $15\,\mathrm{cm}$ 为焦距大小 $10\,\mathrm{cm}$ 的凹透镜。求最终像的位置和横向放大率。
\end{origexamplebox}\par\nopagebreak[3]
\begin{xesolution}
\PrismTrainPanel

先沿图~\ref{ch28:fig:prism-train}把棱镜中的折线光路展开。全反射只把光轴转过 $90^\circ$，不改变会聚性质；棱镜内从上表面经斜面到右表面的轴向路程为 $3+3=6\,\mathrm{cm}$。近轴传播时用约化倾角 $p=n\theta$，高度增量为 $\Delta y=dp/n$，所以该玻璃段的等效空气传播长度是 $6/n=4\,\mathrm{cm}$，不是相位光程 $6n$。故从棱镜右面看，物的等效距离为 $6+4=10\,\mathrm{cm}$；再加棱镜到凸透镜的 $10\,\mathrm{cm}$，有 $u_1=20\,\mathrm{cm}$。

凸透镜：
\[
 v_1=\frac{10\times20}{20-10}=20\,\mathrm{cm},\qquad m_1=-1.
\]
凹透镜位于凸透镜右侧 $15\,\mathrm{cm}$，所以到达它的光原本还要在其右侧 $5\,\mathrm{cm}$ 会聚，即 $u_2=-5\,\mathrm{cm}$。取 $f_2=-10\,\mathrm{cm}$：
\[
 v_2=\frac{(-10)(-5)}{-5-(-10)}=10\,\mathrm{cm},\qquad m_2=-\frac{v_2}{u_2}=2.
\]
因此最终形成在凹透镜右侧
\[
 \boxed{10\,\mathrm{cm}}
\]
处的实像，总横向放大率
\[
 \boxed{m=m_1m_2=-2,\qquad |M|=2}.
\]
负号表示相对棱镜展开后的等效物方向倒立；折回实际空间时还应连同棱镜造成的 $90^\circ$ 方位旋转理解。
\end{xesolution}


\subsection{近轴矩阵方法}


在近轴条件下，用列向量 $(y,p)^{\mathsf T}$ 表示一条光线，其中 $y$ 向上为正，$\theta$ 从向右光轴逆时针量为正，$p=n\theta$ 是约化倾角。这里不用 $u$，以免与物距混淆。
折射率为 $n$ 的均匀介质中传播距离 $d$，以及通过光焦度为 $\Phi$ 的薄透镜，分别对应
\begin{equation}
 \mathbf T(d,n)=
 \begin{pmatrix}1&d/n\\0&1\end{pmatrix},
 \qquad
 \mathbf L(\Phi)=
 \begin{pmatrix}1&0\\-\Phi&1\end{pmatrix}.
\end{equation}
若透镜两侧处于相同折射率 $n$ 的介质中，像方焦距为 $f$，则 $\Phi=n/f$；只有空气中才近似为 $1/f$。各矩阵左乘光线列向量，\textbf{先经过的元件写在最右侧}。例如先传播 $d_1$、过透镜、再传播 $d_2$，总矩阵是 $\mathbf T(d_2,n)\mathbf L(\Phi)\mathbf T(d_1,n)$。一般写为
\begin{equation}
 \begin{pmatrix}y'\\p'\end{pmatrix}
 =\begin{pmatrix}A&B\\C&D\end{pmatrix}
 \begin{pmatrix}y\\p\end{pmatrix},
 \qquad AD-BC=1.
\end{equation}
若固定物平面上的同一点发出不同倾角的光线，在像平面仍交于同一点，则输出高度不能依赖初始倾角，成像条件就是 $B=0$，
横向放大率为 $A$。若平行入射光经过系统后仍保持平行，系统为无焦系统（afocal system），相应条件为 $C=0$。
例如单薄透镜的物平面到像平面矩阵为 $\mathbf T(v,n)\mathbf L(n/f)\mathbf T(u,n)$。乘出右上元为
\[
 B=\frac{u+v-uv/f}{n}.
\]
所以 $B=0$ 恰好恢复 $1/u+1/v=1/f$；左上元 $A=1-v/f=-v/u$ 也恢复带符号倍率。$B$ 的量纲为长度，$C$ 为长度倒数，$A,D$ 无量纲，行列式为一。厚透镜和多透镜组的主平面、焦点与节点，都可从同一矩阵分解得到；这避免了每增加一个折射面就重新建立一套符号法则。



\section{眼睛、视角与放大镜}



把人眼粗略看成空气到折射率 $n=4/3$、曲率半径 $R=5.56\,\mathrm{mm}$ 的单球面系统。远处物体 $u\to\infty$ 时，平行光的像距
\begin{equation}
 f=\frac{nR}{n-1}\approx2.22\,\mathrm{cm},
\end{equation}
与角膜到视网膜的尺度接近。真实人眼还包含角膜、房水、晶状体和玻璃体；水晶体由睫状肌调节曲率，使不同物距的像仍落在固定的视网膜上。水下看不清物体，是因为水的折射率约 $1.33$，与角膜约 $1.376$ 过于接近，角膜折射能力显著下降。

动物眼大体可分为两种工作思路：人眼一类用折射系统在感光面上形成倒立实像；昆虫复眼则由许多小眼单元分别接收不同方向的光。对人眼而言，视网膜位置几乎固定，调焦不能靠前后移动“屏”，只能由睫状肌改变晶状体曲率和等效焦距。看远处时睫状肌较放松；物体移近时，为保持同一像距，系统光焦度必须增大。晶状体形变能力有限，于是出现近点。图~\ref{ch28:fig:s43-20-3} 中角膜、晶状体、视网膜以及光线最终会聚的位置，是理解这一调节链条所必需的四项信息。

\EyePanel

能清楚聚焦的最近点称为近点。青年人的典型明视距离取
\begin{equation}
 L_0=25\,\mathrm{cm}.
\end{equation}
它是目视仪器设计的重要尺度。近点随年龄和屈光状态变化；近视者的远点有限，远视或老视者的近点则较远。$20\,\mathrm{cm}$ 与 $30\,\mathrm{cm}$ 只能作直观示例；实际诊断还需结合远点和屈光度，不能只凭一个距离。

正常眼完全放松时，远处平行光应落在视网膜上。近视眼的系统光焦度相对过强或眼轴偏长，平行光在视网膜前会聚，需要负透镜先使光束稍发散；远视眼或部分老视状态则需正透镜提供额外会聚。显微镜、望远镜的最终像通常设计在无穷远附近，让眼睛放松观察；若使用者近视且不戴镜，调节目镜会把最终虚像移到其有限远点。后面的显微镜例题正是通过“最终虚像应落在哪里”判断镜筒移动方向。



人眼感受到的“大小”由视角决定。尺度 $H$、距离 $L$ 的小物体张角 $\theta\simeq H/L$。光学仪器的角放大率定义为使用仪器后的视角与裸眼基准视角之比。

放大镜是一枚会聚透镜。把物放在焦点以内，透镜形成正立放大的虚像，眼睛靠近透镜观察。若物接近焦点，出射光近似平行，视角 $\theta\simeq H/f$；裸眼把物放在明视距离 $L_0$ 时，$\theta_0\simeq H/L_0$，所以放松眼调焦的角放大率为
\begin{equation}
 M=\frac{\theta}{\theta_0}=\frac{L_0}{f}.
\end{equation}
放大镜的本质，是让眼睛能在小于明视距离的位置清楚观察物体。望远镜、显微镜等目视系统的目镜也起同样作用：它观察前级系统形成的中间像，并进一步放大视角。

若不要求眼睛放松，而把最终虚像调到近点 $v=-L_0$，薄透镜公式给出 $1/u=1/f+1/L_0$。眼睛紧贴透镜时 $\theta\simeq H/u$，相对裸眼近点视角 $H/L_0$ 的角放大率为
\[
 M_{\rm near}=\frac{L_0}{u}=1+\frac{L_0}{f}.
\]
因此教材中常见的 $L_0/f$ 与 $1+L_0/f$ 对应两种不同调焦约定，并不矛盾。前者观察舒适，后者放大略大却更易疲劳。图~\ref{ch28:fig:s43-20-7} 的虚线反向延长线必须交于物体同侧的虚像，不能把放大镜误画成在屏上形成实像。

\MagnifierDiagram



前面的放大镜成像关系还可以统一为目视仪器的基准：这里正常眼明视距离记为 $D_0=L_0$，以免与口径 $D$ 混淆。光学仪器的“放大”通常指视角增大，而不是在屏上形成更大的实像。放大镜把物体放在焦点附近，
使眼睛在较舒适状态下观察放大的虚像。最终像在无穷远时角放大率
\begin{equation}
 M_\infty=\frac {D_0}f;
\end{equation}
最终像在近点时
\begin{equation}
 M_{D_0}=1+\frac {D_0}f.
\end{equation}
两式对应不同调焦条件，不能混用。



\subsection{角放大率的比较对象}
放大镜产生正立放大的虚像，但“虚像高变大了多少”并不是眼睛感到放大多少的完整答案。眼睛接收的是方向范围，即物体或像对眼睛张开的视角。比较时应固定同一物体，并明确裸眼基准通常取明视距离 $D$，于是裸眼视角为 $\alpha_0\simeq h/D$。

眼睛紧靠薄透镜、最终像在无穷远时，物位于前焦面，出射束视角为 $\alpha\simeq h/f$，所以 $M=\alpha/\alpha_0=D/f$。眼睛放松与物在焦面是同一个光路条件的两种表述，不能同时把物距取为任意值。

若最终虚像在距离透镜约 $D$ 的近点，取 $v=-D$，成像公式给出
\[
 \frac1u=\frac1f+\frac1D.
\]
眼睛靠近透镜时 $\alpha\simeq h/u$，因此 $M=D/u=1+D/f$。两个常用公式不是互相矛盾的近似，而是对应不同调焦状态。眼睛与透镜间距离不可忽略时，近点应从眼睛量起，须先把这一距离加入几何关系后再求视角。

视力分辨、角放大率和视场也不是一回事。放大率可以增加，但像的细节若已被衍射或像差模糊，继续放大只会放大模糊斑。有限镜片口径还会限制能同时看到的物体范围；判断视场需要画边缘光线，不能只看过光心的一条线。


\section{显微镜与望远镜}

\subsection{显微镜的中间像与调焦}


复合显微镜由短焦物镜和目镜组成。物镜先形成放大的倒立实像，目镜再把中间像作为物体形成放大虚像。图\ref{ch28:fig:microscope}中，光学筒长 $L$ 定义为物镜后焦点 $F'_o$ 到目镜前焦面（中间像平面）的轴向距离，不是两个镜片之间的机械距离。物镜像距为 $v_o=f_o+L$，从薄透镜公式有 $m_o=1-v_o/f_o=-L/f_o$。目镜角放大率为 $D_0/f_e$，因此最终像在无穷远的近轴角放大率为
\begin{equation}
 M\simeq-\frac{L}{f_o}\frac{D_0}{f_e}.
\end{equation}
负号表示最终像相对原物倒立。提高几何放大率不能突破衍射分辨极限；有用放大应与物镜数值孔径相匹配。

\begin{center}\begin{minipage}{\linewidth}\centering
\begin{tikzpicture}[>=Latex,font=\small,every node/.style={fill=none}]
\draw[gray,thin,->] (-2.0000000,0.0000000)--(8.2000000,0.0000000);
\draw[DeepBlue,thick,<->] (0.0000000,-2.1000000)--(0.0000000,2.1000000);
\draw[DeepBlue,thick,<->] (6.2000000,-2.1000000)--(6.2000000,2.1000000);
\node[above,fill=none,inner sep=2pt] at (0.0000000,2.2000000) {物镜};
\node[above,fill=none,inner sep=2pt] at (6.2000000,2.2000000) {目镜};
\draw[black,thick,->] (-1.2500000,0.0000000)--(-1.2500000,0.3000000);
\fill (5.0000000,-1.2000000) circle (1.2pt);
\node[below,fill=none,inner sep=2pt] at (5.0000000,-1.2000000) {};
\draw[gray,dashed,thin] (5.0000000,-1.2000000)--(4.5000000,1.1500000);
\node[above,fill=none,inner sep=2pt] at (4.5000000,1.2000000) {中间实像};
\draw[gray,dashed,thin] (5.0000000,-1.5000000)--(5.0000000,1.4000000);
\draw[DeepBlue,thick,->] (-1.2500000,0.3000000)--(0.0000000,0.6500000);
\draw[DeepBlue,thick,->] (0.0000000,0.6500000)--(6.2000000,-1.6440000);
\draw[DeepBlue,thick,->] (6.2000000,-1.6440000)--(8.0000000,0.1560000);
\draw[DeepBlue,thick,->] (-1.2500000,0.3000000)--(0.0000000,-0.1500000);
\draw[DeepBlue,thick,->] (0.0000000,-0.1500000)--(6.2000000,-1.4520000);
\draw[DeepBlue,thick,->] (6.2000000,-1.4520000)--(8.0000000,0.3480000);
\draw[gray,<->] (1.0000000,-1.8500000)--(5.0000000,-1.8500000);
\node[below,fill=none,inner sep=2pt] at (3.0000000,-1.8500000) {$L$};
\fill (1.0000000,0.0000000) circle (1.2pt);
\node[below,fill=none,inner sep=2pt] at (1.0000000,0.0000000) {};
\draw[gray,dashed,thin] (1.0000000,0.0000000)--(1.0000000,-0.9500000);
\node[below,fill=none,inner sep=2pt] at (1.0000000,-1.0000000) {$F'_o$};
\draw[gray,<->] (5.0000000,-2.5000000)--(6.2000000,-2.5000000);
\node[below,fill=none,inner sep=2pt] at (5.6000000,-2.5000000) {$f_e$};
\end{tikzpicture}
\figcap{显微镜的中间像在目镜前焦面；两透镜之间不发生折转。$L$ 从物镜后焦点量到中间像平面。\label{ch28:fig:microscope}}
\end{minipage}\end{center}



\Needspace{8\baselineskip}\begin{origexamplebox}{【小明逞强】}\label{ch28:ex:s43-20-4}
近视的小明不戴眼镜，与视力正常的小红使用同一显微镜。小红调清楚后，小明应上调还是下调镜筒才能看清？
\end{origexamplebox}\par\nopagebreak[3]
\begin{xesolution}
近视眼不戴镜时远点较近，需要显微镜给出较近的最终虚像。目镜焦距为 $f_e$，中间像在目镜焦内、物距为 $u_e$，虚像距离大小 $|v_e|=f_eu_e/(f_e-u_e)$；要使 $|v_e|$ 变小，应使 $u_e$ 变小。

下调镜筒使物镜靠近标本。对物镜而言，物距更接近焦距，所成中间实像距物镜变大；目镜与物镜间距固定，因此中间像到目镜的距离 $u_e$ 变小，最终虚像变近。故应
\[
 \boxed{\text{下调镜筒（使物镜略靠近标本）}}.
\]
把调节方向写成导数也能自检。镜筒向标本移动 $\varepsilon>0$ 时，物镜物距变化 $\Delta u_o=-\varepsilon$，而
\[
 \frac{\dd v_o}{\dd u_o}=-\frac{f_o^2}{(u_o-f_o)^2}<0,
 \qquad u_e=\ell-v_o,
\]
其中 $\ell$ 为固定的两镜间距。因此 $v_o$ 增大而 $u_e$ 减小。若小红原来按放松眼调到无穷远，而小明远点距目镜为 $F_R$（忽略眼距），所需 $u_e=f_eF_R/(F_R+f_e)<f_e$，与上述方向一致。这里比较两人都采用舒适调焦的状态；若眼睛主动调节或佩戴矫正镜，应重新指定最终像的位置。
实际操作应缓慢并观察侧面距离，避免物镜碰到玻片。
\end{xesolution}
\subsection{望远镜的无焦系统与有效倍率}


开普勒望远镜（Keplerian telescope）用正焦距物镜和正焦距目镜。远物经物镜在焦平面形成实像，
调节为无焦系统时两透镜焦平面重合，角放大率
\begin{equation}
 M=-\frac{f_o}{f_e},
\end{equation}
镜筒长约 $f_o+f_e$。图\ref{ch28:fig:telescope}中，入射倾角 $\theta$ 对应中间像高 $f_o\theta$；此像在目镜前焦面，所以出射倾角为 $-f_o\theta/f_e$，两条出射光平行。伽利略望远镜（Galilean telescope）使用负焦距目镜，成正立虚像，镜筒更短，但视场受限。

\begin{center}\begin{minipage}{\linewidth}\centering
\begin{tikzpicture}[>=Latex,font=\small,every node/.style={fill=none}]
\draw[gray,thin,->] (-2.0000000,0.0000000)--(7.5000000,0.0000000);
\draw[DeepBlue,thick,<->] (0.0000000,-2.1000000)--(0.0000000,2.1000000);
\draw[DeepBlue,thick,<->] (5.5000000,-2.1000000)--(5.5000000,2.1000000);
\node[above,fill=none,inner sep=2pt] at (0.0000000,2.2000000) {物镜};
\node[above,fill=none,inner sep=2pt] at (5.5000000,2.2000000) {目镜};
\fill (4.0000000,-0.4800000) circle (1.2pt);
\node[below,fill=none,inner sep=2pt] at (4.0000000,-0.4800000) {};
\draw[gray,dashed,thin] (4.0000000,-0.4800000)--(3.5000000,1.1500000);
\node[above,fill=none,inner sep=2pt] at (3.5000000,1.2000000) {中间像点};
\draw[gray,dashed,thin] (4.0000000,-1.5000000)--(4.0000000,1.4000000);
\draw[DeepBlue,thick,->] (-1.6000000,0.8420000)--(0.0000000,0.6500000);
\draw[DeepBlue,thick,->] (0.0000000,0.6500000)--(5.5000000,-0.9037500);
\draw[DeepBlue,thick,->] (5.5000000,-0.9037500)--(7.3000000,-0.3277500);
\draw[DeepBlue,thick,->] (-1.6000000,0.0420000)--(0.0000000,-0.1500000);
\draw[DeepBlue,thick,->] (0.0000000,-0.1500000)--(5.5000000,-0.6037500);
\draw[DeepBlue,thick,->] (5.5000000,-0.6037500)--(7.3000000,-0.0277500);
\draw[gray,<->] (4.0000000,-2.5000000)--(5.5000000,-2.5000000);
\node[below,fill=none,inner sep=2pt] at (4.7500000,-2.5000000) {$f_e$};
\end{tikzpicture}
\figcap{无焦望远镜：公共焦面处光线相交但不折转；目镜后的两条出射线严格平行。\label{ch28:fig:telescope}}
\end{minipage}\end{center}

望远镜分辨本领取决于物镜口径而非焦距。圆孔的瑞利判据（Rayleigh criterion）为
\begin{equation}
 \theta_{\min}=1.22\frac\lambda D.
\end{equation}
增大口径既提高集光能力，也提高角分辨率；仅增大目镜倍率不会恢复物镜没有记录的细节。

\Needspace{8\baselineskip}\begin{origexamplebox}{望远镜口径的分辨要求}
两颗星的角距离为 $0.50''$，用波长 $550\,\mathrm{nm}$ 的光观察。按瑞利判据估算望远镜最小口径。
\end{origexamplebox}\par\nopagebreak[3]
\begin{xesolution}
先把角秒换成弧度：
\[
 \theta=0.50\frac{\pi}{180\times3600}
 \approx2.42\times10^{-6}\,\mathrm{rad}.
\]
所以
\[
 D_{\min}=1.22\frac\lambda\theta
 \approx1.22\frac{550\times10^{-9}}{2.42\times10^{-6}}
 \approx0.277\,\mathrm m.
\]
即 $\boxed{D_{\min}\approx28\,\mathrm{cm}}$。真实观测还会受大气视宁度、像差和信噪比限制。
\end{xesolution}



\Needspace{8\baselineskip}\begin{origexamplebox}{【天文望远镜】}\label{ch28:ex:s43-20-5}
天文台常见双凸式透射望远镜。（1）为什么不采用凹凸式（伽利略式）作主要天文望远镜？（2）怎样简易判断望远镜的性能高低？
\end{origexamplebox}\par\nopagebreak[3]
\begin{xesolution}
双凸的开普勒结构由正物镜先形成实的中间像，再由正目镜放大。它虽然给出倒像，但视场较大，容易在中间焦面放置十字丝、相机、光谱狭缝等。凹凸伽利略结构成正像且镜筒短，却没有可接收的中间实像，出瞳位于系统内部，视场较小，不适合主要天文观测附件。

简易比较不能只看标称“倍数”。在大气视宁度和加工质量相当时，物镜有效口径 $D$ 决定集光能力 $\propto D^2$，并决定衍射角分辨率 $\theta\simeq1.22\lambda/D$；因此先看\textbf{有效口径和像质}，再看焦距、消色差、镀膜、支架稳定性和目镜配置。超过口径所允许的倍率只会把模糊斑放大。
\end{xesolution}

\section{仪器分辨本领与光能探测}



棱镜或光栅把不同波长分到不同方向，准直镜把入射狭缝变成平行光，成像镜把各方向重新聚焦到探测器。
平面光栅正入射时 $d\sin\theta_m=m\lambda$。若 $N$ 条刻线均匀相干照明，忽略像差，在非零级次 $m$ 的主极大附近，用瑞利判据得到理论分辨本领
\begin{equation}
 \mathcal R=\frac\lambda{\Delta\lambda}=|m|N.
\end{equation}
这是因为两波长的主极大角间距为 $|m|\Delta\lambda/(d\cos\theta_m)$，单条谱线主极大到第一极小的角距为 $\lambda/(Nd\cos\theta_m)$；令二者相等即可消去共同几何因子。$m=0$ 没有通常意义的角分光能力，不能作为上述分辨公式的工作级次；接近掠射时也需谨慎使用局部小角展开。狭缝宽度、像差和探测器采样会降低实际分辨率；狭缝无限变窄又会显著损失光通量。



辐射度学（radiometry）按电磁辐射携带的功率计量；光度学（photometry）还要按明视觉光谱响应 $V(\lambda)$ 加权。
若光谱辐射通量为 $\Phi_{e,\lambda}$，光通量为
\begin{equation}
 \Phi_v=K_m\int_0^\infty V(\lambda)\Phi_{e,\lambda}\,\dd\lambda,
 \qquad K_m\simeq683\,\mathrm{lm/W}.
\end{equation}
国际单位制用频率 $540\,\mathrm{THz}$ 的单色辐射把光视效能固定为 $683\,\mathrm{lm/W}$；其他波长即使辐射功率相同，
视觉亮度也会因 $V(\lambda)$ 不同而改变。常用光度量为
\begin{equation}
 I_v=\frac{\dd\Phi_v}{\dd\Omega},\qquad
 E_v=\frac{\dd\Phi_v}{\dd A},\qquad
 L_v=\frac{\dd^2\Phi_v}{\dd A\cos\theta\,\dd\Omega},
\end{equation}
分别称发光强度、照度和亮度，单位依次为坎德拉、勒克斯和坎德拉每平方米。光电二极管、CCD/CMOS 等探测器有各自的
光谱响应曲线；把“人眼看起来一样亮”直接当成“探测器接收功率相同”通常是错误的。

\tcbset{breakable/.default=true}\begin{starread}{三刺激值与 CIE 1931 色度图}
人眼明视觉主要由三类视锥细胞提供长、中、短波段响应。CIE 1931 标准把光谱 $P(\lambda)$ 映射为三刺激值
\[
 X=\int\bar x(\lambda)P(\lambda)\,\dd\lambda,\quad
 Y=\int\bar y(\lambda)P(\lambda)\,\dd\lambda,\quad
 Z=\int\bar z(\lambda)P(\lambda)\,\dd\lambda,
\]
再用 $x=X/(X+Y+Z)$、$y=Y/(X+Y+Z)$ 表示与总亮度分离的色度。两束光相加时，三刺激值线性相加；
在图\ref{ch28:fig:chromaticity}的色度平面上，二者混合色位于两点连线上。显示器选定三原色后，只能覆盖三原色三角形内部的颜色，并不能覆盖整个可见色域。

\begin{center}\begin{minipage}{\linewidth}\centering
\begin{tikzpicture}
\begin{axis}[width=.70\linewidth,height=6.0cm,axis lines=left,
  xlabel={$x$},ylabel={$y$},xmin=0,xmax=.80,ymin=0,ymax=.88,
  xtick={0,.2,.4,.6,.8},ytick={0,.2,.4,.6,.8},
  every axis label/.style={fill=none}]
  \addplot[curve,very thick] coordinates{
    (.174,.005)(.157,.018)(.124,.058)(.091,.133)(.045,.295)(.008,.538)
    (.014,.750)(.074,.834)(.155,.806)(.230,.754)(.302,.692)(.373,.625)
    (.444,.555)(.513,.487)(.575,.424)(.627,.373)(.666,.334)(.692,.308)
    (.708,.292)(.719,.281)(.726,.274)(.735,.265)(.174,.005)};
  \addplot[DeepBlue,thick] coordinates{(.64,.33)(.30,.60)(.15,.06)(.64,.33)};
  \addplot[only marks,mark=*,mark size=1.7pt] coordinates{(.333,.333)};
  \node[anchor=west] at (axis cs:.35,.35) {$E$};
  \node[DeepBlue] at (axis cs:.52,.10) {RGB 色域};
  \draw[->,thin] (axis cs:.49,.13)--(axis cs:.43,.30);
  \node at (axis cs:.20,.66) {可见色域};
\end{axis}
\end{tikzpicture}
\figcap{CIE 1931 色度示意：弯曲边界为近似光谱轨迹，两端的直线为紫线（非单色光）。E 为等能光谱白点 $(1/3,1/3)$，不等同于任意 RGB 标准的白点。\label{ch28:fig:chromaticity}}
\end{minipage}\end{center}
\end{starread}\tcbset{breakable/.default=false}

\tcbset{breakable/.default=true}\begin{starread}{光阑、光瞳、像差与展度守恒}
孔径光阑限制轴上光束锥角，视场光阑限制可见物方范围。孔径光阑经其前方光具组成的像称入射光瞳，经其后方光具组成的像称出射光瞳；
望远镜目镜应让出射光瞳与人眼瞳孔匹配。球差、彗差、像散、场曲、畸变和色差具有不同的孔径与视场依赖，
不能只用“焦点不准”概括。

无损系统中的微分展度
\[
 \dd G=n^2\cos\theta\,\dd A\,\dd\Omega
\]
沿光路守恒。这意味着被动光学系统不能同时把光束横截面积和发散立体角任意压缩；提高照度、扩大视场、增大口径和保持像质之间
必然存在约束，而不是只要增加透镜片数就能全部兼得。
\end{starread}\tcbset{breakable/.default=false}

\endgroup
\subsection{将放大、分辨和亮度分开判断}
用望远镜观测两个角距为 $\Delta\theta$ 的远处点源时，物镜首先把它们形成相邻的衍射斑。口径为 $D_o$ 的圆形物镜，其瑞利判据为 $\Delta\theta\gtrsim1.22\lambda/D_o$。目镜只是把已经形成的像以更大的视角交给眼睛；若两斑在物镜焦面上不能分开，仅更换更短焦距的目镜不能恢复丢失的细节。

相机还需比较衍射斑和像素。焦距为 $f$ 时，艾里斑第一暗环半径约为 $1.22\lambda f/D_o$。减小孔径可以改善部分几何像差和增加景深，却会增大这个衍射尺度。由此可见“缩光圈越小越清晰”并非普遍规律，清晰度取决于像差、衍射、采样和对焦误差中哪一项占主导。

最后检查能量量纲：光功率 $\Phi_e$ 的单位为瓦，照度型的辐照度 $E_e=\dd\Phi_e/\dd A$ 的单位为瓦每平方米，曝光量 $H_e=\int E_e\dd t$ 的单位为焦耳每平方米。像变大并不创造能量；在接收总功率近似不变时，同样的功率分布到更大像面积上，单位面积得到的能量反而减少。光度量还经过人眼光谱响应加权，不能把流明与瓦不加条件地互换。

